\ifdefined\SubmissionReview
\documentclass[11pt]{article}
\else
\documentclass[10pt,twocolumn]{article}
\fi
\usepackage[a4paper,top=18mm,bottom=20mm,left=17mm,right=17mm,columnsep=7mm]{geometry}
\usepackage{times,amsmath,amssymb,amsthm,mathtools,bm,booktabs,multirow,array,graphicx,xcolor,url,enumitem,float,balance,subcaption}
\usepackage[font=small,labelfont=bf]{caption}
\usepackage{placeins}
\usepackage{cuted}
\usepackage{tikz}
\usetikzlibrary{arrows,positioning}
\newenvironment{widetable}
  {\begin{strip}\centering\captionsetup{type=table}}
  {\end{strip}}
\setlength{\parindent}{1em}\setlength{\parskip}{0pt}
\raggedbottom
\setcounter{topnumber}{5}
\setcounter{bottomnumber}{3}
\setcounter{totalnumber}{8}
\setcounter{dbltopnumber}{5}
\renewcommand{\topfraction}{0.95}
\renewcommand{\bottomfraction}{0.90}
\renewcommand{\textfraction}{0.05}
\renewcommand{\floatpagefraction}{0.72}
\renewcommand{\dbltopfraction}{0.95}
\renewcommand{\dblfloatpagefraction}{0.72}
\setlength{\textfloatsep}{12pt plus 2pt minus 2pt}
\setlength{\floatsep}{11pt plus 2pt minus 2pt}
\makeatletter
\setlength{\@fptop}{0pt}
\setlength{\@fpsep}{11pt}
\setlength{\@fpbot}{0pt plus 1fil}
\setlength{\@dblfptop}{0pt}
\setlength{\@dblfpsep}{11pt}
\setlength{\@dblfpbot}{0pt plus 1fil}
\makeatother
\allowdisplaybreaks

\newtheorem{theorem}{Theorem}
\newtheorem{proposition}{Proposition}
\newtheorem{corollary}{Corollary}
\theoremstyle{remark}
\newtheorem{remark}{Remark}

\newcommand{\R}{\mathbb{R}}
\newcommand{\Tr}{\operatorname{Tr}}
\newcommand{\eps}{\varepsilon}
\newcommand{\ones}{\mathbf{1}}
\newcommand{\argmaxop}{\operatorname*{arg\,max}}
\newcommand{\had}{\odot}

\title{Evidence-guided conservative graph calibration for direct semi-supervised clustering}
\author{Deshu Sun, Feng Wang$^{*}$\\
\small College of Data Science and Information Engineering, Guizhou Minzu University, Guiyang, Guizhou 550025, P.R. China\\
\small $^{*}$Corresponding author: wangfeng@gzmu.edu.cn\\
\small Manuscript prepared for Pattern Analysis and Applications}
\date{}
\ifdefined\SubmissionReview
\newcommand{\frontblock}[1]{#1}
\else
\newcommand{\frontblock}[1]{\twocolumn[\begin{@twocolumnfalse}#1\end{@twocolumnfalse}]}
\fi

\begin{document}
\frontblock{
\maketitle
\begin{abstract}
Graph-based semi-supervised clustering usually relies on a prescribed neighborhood graph, although limited labeled data can also provide task-specific evidence about whether that graph should be modified. We formulate graph modification as an evidence-guided Conservative Graph Calibration (CGC) problem for Laplacian-regularized semi-supervised representation models. A binary baseline graph and a same-support locally weighted candidate are evaluated by out-of-fold harmonic label propagation. Their foldwise predictive-loss difference estimates the advantage of local weighting, while its foldwise variability discounts unstable evidence. The resulting retained-advantage coefficient determines how far the final graph moves toward the candidate and yields exact baseline recovery when the evidence for modification is insufficient. CGC is completed before downstream optimization and introduces neither a graph-learning variable nor an additional objective coefficient. We first instantiate CGC in one-hot-constrained NMF. Across six benchmarks, CGC-GOCNMF improves its parent on COIL20, COIL100, Optdigits, and MNIST while retaining essentially baseline behavior on PIE and YaleB. We then transfer the same calibration rule, without retuning it, to GNMFLD, whose supervision uses a label-discrimination penalty rather than fixed one-hot rows. CGC-GNMFLD improves ACC and NMI in all 20 paired trials on each benchmark, with mean ACC gains of $0.24$-$4.57$ percentage points and NMI gains of $0.57$-$3.06$ points. Industrial-image, representation-robustness, ablation, sensitivity, and controlled-corruption experiments further characterize when candidate weighting is retained, attenuated, or rejected. These results support CGC as a reusable pre-optimization mechanism for regulating structural graph modification from limited label evidence.
\end{abstract}
\vspace{1mm}\noindent\textbf{Keywords:} direct semi-supervised clustering; nonnegative matrix factorization; conservative graph calibration; label evidence; graph regularization; variability-adjusted shrinkage; steel surface defect clustering
\vspace{5mm}
}

\section{Introduction}
\label{sec:introduction}

Nonnegative matrix factorization (NMF) decomposes a nonnegative data matrix into two low-rank nonnegative factors and provides an interpretable representation of high-dimensional data \cite{Lee1999NMF,Lee2000Algorithms}. Owing to its parts-based representation and simple optimization structure, NMF has been widely used for dimensionality reduction, image representation, and clustering. Standard NMF, however, is unsupervised and does not directly exploit either the local geometry of the data or the labels available for a small subset of samples.

Graph-regularized NMF (GNMF) incorporates manifold information by encouraging neighboring samples to have similar low-dimensional representations \cite{Cai2011GNMF}. Semi-supervised NMF further uses partial labels through hard constraints, label-discrimination losses, or propagated supervisory information \cite{Liu2012CNMF,Mo2024PNLPSCNMF,Xing2021GNMFLD}. More elaborate approaches learn multiple graph weights, update neighborhood structures jointly with the factors, or combine NMF with label propagation and hypergraph learning \cite{Jing2026SNMFWLP,Li2023ABNMTF,Ren2025AHSSNMF,Yang2025SPNMF,Yi2020AGSSNMF,Zhang2021MSNMF}. These methods can improve clustering, but they often introduce additional variables, regularization coefficients, or iterative subproblems. Many NMF clustering methods also require $K$-means or another post-processing procedure, so the final assignments are not obtained directly from the factorization.

Graph regularized one-hot-constrained NMF (GOCNMF) provides a concise alternative for direct semi-supervised clustering \cite{Li2025GOCNMF}. It fixes the representations of labeled samples as one-hot vectors and uses graph regularization to guide the unlabeled representations toward class-indicative coordinates. The cluster assignment of each sample is therefore obtained from the largest coordinate of its representation, without an external classifier or $K$-means. This architecture preserves the simplicity of NMF and avoids an additional label-fitting coefficient.

The reliability of graph regularization nevertheless depends on the prescribed neighborhood structure. A binary $p$-nearest-neighbor graph is simple and stable, but it ignores variations in local similarity strength. A locally scaled graph can encode such variations more faithfully \cite{Zelnik2004SelfTuning}, yet greater expressiveness does not necessarily imply greater task relevance. In a semi-supervised setting, the available labels provide direct evidence about whether changing the graph improves class-consistent propagation. This leads to a distinct methodological problem: graph modification should itself be treated as a decision under limited supervisory evidence rather than as an unconditional preprocessing choice.

This perspective assigns the labeled samples two complementary roles. First, they supervise the downstream representation according to the parent model, for example through fixed one-hot rows or a label-discrimination loss. Second, they provide task-specific evidence for deciding how far the graph should depart from a trusted baseline. The second role is usually implicit in fixed-graph methods or absorbed into a larger joint graph-learning model. Here it is treated explicitly through a pre-factorization calibration mechanism. We therefore study whether limited labeled evidence can regulate structural graph modification before representation learning, and whether instability in the observed graph advantage can be used to control the amount of deviation from the baseline.

To address this problem, we develop an evidence-guided Conservative Graph Calibration (CGC) mechanism. CGC compares a prescribed baseline graph $W^{(0)}$ with a same-support candidate graph $W^{(1)}$ using out-of-fold harmonic label propagation \cite{Zhu2003Harmonic} and multiclass squared prediction loss \cite{Brier1950}. If $\Delta$ denotes the mean foldwise advantage of the candidate graph and $s$ its foldwise variability, the retained-advantage coefficient is
\[
\theta=
\begin{cases}
\dfrac{(\Delta-s)_{+}}{\Delta},&\Delta>0,\\[1mm]
0,&\Delta\le0,
\end{cases}
\]
with $W_{\theta}=(1-\theta)W^{(0)}+\theta W^{(1)}$.
Thus, only the fraction of the observed advantage that survives the variability discount is allowed to modify the graph. When $\Delta\le s$, the coefficient is zero and the baseline is recovered exactly.

We first instantiate CGC in GOCNMF, yielding CGC-GOCNMF. Calibration is completed before factorization and is not optimized jointly with the factor matrices, so the fixed one-hot constraints, multiplicative solver, and direct row-wise $\argmaxop$ assignment rule of the parent are retained. To test whether the same calibration mechanism depends on this hard one-hot supervision, we further transfer CGC to GNMFLD \cite{Xing2021GNMFLD}, which instead uses a label-discrimination penalty. In both instantiations, CGC changes only the graph supplied to the parent optimizer.

The main contributions are summarized as follows.

\begin{enumerate}[leftmargin=2em]
\item \textbf{Evidence-guided conservative graph calibration.} We formulate graph modification as a pre-optimization decision for Laplacian-regularized semi-supervised representation models. A baseline graph and a same-support candidate are evaluated from out-of-fold labeled predictions rather than unconditionally replacing or jointly learning the graph.

\item \textbf{Variability-adjusted retained advantage.} CGC discounts a positive mean graph advantage by its fold-to-fold variability. The resulting coefficient has an explicit retained-advantage interpretation, is monotone in the evidence and variability terms, and recovers the baseline exactly when the evidence for graph modification is insufficient.

\item \textbf{Controlled structural perturbation and parent-model preservation.} We show that the calibration coefficient scales the graph and Laplacian deviations exactly and bounds the induced change in any downstream objective of the stated Laplacian-regularized form. For the primary GOCNMF instantiation, we additionally establish graph validity, monotonic objective values under the exact multiplicative updates, and stationarity of strictly positive fixed points.

\item \textbf{Two-parent empirical validation.} The primary CGC-GOCNMF study combines six benchmarks, an industrial image dataset, representation robustness, ablation, sensitivity analysis, controlled graph corruption, and timing. A second controlled transfer to GNMFLD uses the same calibration rule with no retuning and tests whether CGC remains effective under a distinct label-discrimination supervision mechanism.
\end{enumerate}

The remainder of this paper is organized as follows. Section~\ref{sec:related} reviews related semi-supervised and graph-learning NMF methods. Section~\ref{sec:method} introduces the candidate graphs, out-of-fold assessment, and conservative calibration rule. Section~\ref{sec:analysis} presents the structural, convergence, and complexity analyses. Section~\ref{sec:experiments} reports the experimental results. Section~\ref{sec:conclusion} concludes the paper.

\section{Related work}
\label{sec:related}

\subsection{Semi-supervised NMF}

Semi-supervised NMF uses limited labels through hard constraints, discriminative losses, propagated supervision, or self-generated supervisory structure. Representative examples include CNMF \cite{Liu2012CNMF}, GNMFLD \cite{Xing2021GNMFLD}, robust and dual-supervision variants \cite{Peng2021RSNMF,Peng2022DCNMF}, self-supervised semi-supervised NMF \cite{Chavoshinejad2023S4NMF}, EDNMF \cite{Li2025EDNMF}, and methods based on label or pairwise-constraint propagation \cite{Mo2024PNLPSCNMF,Jing2026SNMFWLP,Yang2025SPNMF}. Deep semi-supervised graph clustering further combines nonlinear representation learning with graph-based supervision \cite{Qi2024GSSC}.

These methods primarily change how supervision enters representation learning. CGC addresses the complementary structural question of whether a prescribed graph should be modified before the downstream model is optimized. Harmonic propagation is used only to compare candidate graphs from held-out labeled scores; no pseudo-label matrix is inserted into the downstream objective.

\subsection{Multiple and adaptive graph learning}

GNMF introduced Laplacian regularization to preserve local manifold geometry \cite{Cai2011GNMF}. Later work learned multiple-graph weights, adaptive neighbors, or higher-order structures jointly with the representation, including MSNMF \cite{Zhang2021MSNMF}, AGSSNMF \cite{Yi2020AGSSNMF}, ABNMTF \cite{Li2023ABNMTF}, and hypergraph variants \cite{Ren2025AHSSNMF,Yin2023HSSNMF}. These approaches optimize graph-related quantities inside the learning model.

CGC instead compares two fixed same-support graphs before downstream optimization. Held-out label-prediction quality determines the direction of graph modification, while foldwise variability contracts unstable evidence. This differs from hard validation-based graph selection, grid search over mixture weights, and joint adaptive-graph learning. It is also distinct from safety-aware graph-based semi-supervised learning, where graph selection is coupled with a safety-oriented learner \cite{Gan2018SaGSSL}: CGC regulates how far a prescribed graph is allowed to move before the parent representation model is optimized. The variability term is therefore used as a deterministic stability discount rather than as a statistical confidence bound.

\begin{widetable}
\centering
\caption{Methodological positioning of representative semi-supervised NMF methods.}
\label{tab:method_comparison}
\footnotesize
\setlength{\tabcolsep}{16.0pt}
\begin{tabular}{llllcc}
\toprule
Method & Label use & Graph handling & Decision stage & Direct & Recovery \\
\midrule
GNMFLD & Discrimination loss & Fixed graph & Pre-factorization & Yes & -- \\
CGC-GNMFLD & Same + OOF evidence & Calibrated fixed pair & Pre-factorization & Yes & Exact \\
EDNMF & Element constraints & No graph & -- & Yes & -- \\
MSNMF & Pairwise constraints & Adaptive multi-graph & Joint & No & No \\
ABNMTF & Label block structure & Adaptive neighbors & Joint & No & No \\
GOCNMF & Fixed one-hot rows & Binary $p$-NN & Pre-factorization & Yes & -- \\
CGC-GOCNMF & One-hot + OOF evidence & Calibrated fixed pair & Pre-factorization & Yes & Exact \\
\bottomrule
\end{tabular}%
\end{widetable}

GNMFLD and GOCNMF use fixed graphs, whereas MSNMF and ABNMTF learn graph-related quantities during optimization. The two CGC instantiations retain their respective parent supervision mechanisms and replace only the prescribed graph by an evidence-calibrated convex combination of two fixed same-support candidates; when retained evidence vanishes, the parent graph is recovered exactly.

\subsection{Direct clustering with one-hot constraints}

Many NMF-based clustering methods require $K$-means or another post-processing procedure to obtain the final assignments. GOCNMF instead fixes the labeled representations as one-hot vectors and directly assigns each unlabeled sample according to its largest representation coordinate \cite{Li2025GOCNMF}. This avoids an external classifier and an additional label-fitting coefficient.

CGC-GOCNMF preserves the factorization structure, one-hot constraints, and direct assignment rule of GOCNMF while introducing a separate evidence-guided decision for graph modification before factorization. The transfer instantiation CGC-GNMFLD preserves GNMFLD's label-discrimination term and direct class-coordinate assignment instead. In both cases, the calibration coefficient is determined from out-of-fold labeled predictions before factorization; when $\Delta\le s$, it becomes zero and the original parent graph is recovered exactly.

\section{Evidence-guided conservative graph calibration}
\label{sec:method}

\subsection{General calibration setting}

Let
\[
X=[x_1,\ldots,x_n]\in\R_{+}^{m\times n}
\]
be a nonnegative data matrix whose columns are normalized to unit Euclidean norm. The samples belong to $c$ classes. Let $\mathcal L$ and $\mathcal U$ denote the labeled and unlabeled index sets, respectively, where $\mathcal L\cap\mathcal U=\varnothing$ and $\mathcal L\cup\mathcal U=\{1,\ldots,n\}$. For $i\in\mathcal L$, let $y_i\in\{1,\ldots,c\}$ be the class label and $e_{y_i}\in\{0,1\}^{c}$ the corresponding one-hot vector. Denote the matrix of labeled one-hot vectors by
\[
C_{\mathcal L}=\bigl[e_{y_i}^{T}\bigr]_{i\in\mathcal L}.
\]

CGC is defined before the downstream representation model is optimized. It takes a baseline graph $W^{(0)}$, a candidate graph $W^{(1)}$, and the labeled subset as input; it then uses held-out label evidence to construct a calibrated graph $W_\theta$. To make the scope of this mechanism explicit, consider a generic graph-regularized objective
\begin{equation}
\label{eq:generic_graph_objective}
\mathcal J_L(Z)
=
\Phi(Z)
+
\gamma\Tr\!\bigl(H(Z)^T L H(Z)\bigr),
\qquad \gamma>0,
\end{equation}
where $Z$ denotes the model variables, $\Phi$ contains the graph-independent part of the objective, and $H(Z)\in\R^{n\times r}$ is the representation on which graph regularization acts. CGC modifies only the Laplacian supplied to \eqref{eq:generic_graph_objective}; it does not alter $\Phi$, introduce a graph variable into $Z$, or prescribe a particular downstream solver. Graph-regularized NMF is a direct instance with $H(Z)=V$. Sections~\ref{subsec:goc_instantiation} and~\ref{subsec:gnmfld_instantiation} give two downstream instantiations with different supervision mechanisms.

\subsection{Binary and locally scaled graphs}

For the normalized samples, define the cosine distance
\begin{equation}
\label{eq:distance}
d_{ij}=\max\{0,1-x_i^{T}x_j\}.
\end{equation}
Let $\mathcal N_p(i)\subseteq\{1,\ldots,n\}\setminus\{i\}$ contain the $p$ nearest neighbors of sample $i$, and define the undirected support
\[
\mathcal E_p=\bigl\{\{i,j\}:1\le i<j\le n,\ j\in\mathcal N_p(i)\ \text{or}\ i\in\mathcal N_p(j)\bigr\}.
\]
The binary graph is
\begin{equation}
\label{eq:binary_graph}
w_{ij}^{(0)}=
\begin{cases}
1,&\{i,j\}\in\mathcal E_p,\\
0,&\text{otherwise},
\end{cases}
\qquad
w_{ii}^{(0)}=0.
\end{equation}

Let $i_p$ be the $p$-th nearest neighbor of $i$, and set
\begin{equation}
\label{eq:local_scale}
\sigma_i=\max\{d_{i,i_p},\eps\}.
\end{equation}
Following the local-scaling principle \cite{Zelnik2004SelfTuning}, define
\begin{equation}
\label{eq:directed_local_graph}
a_{ij}=
\begin{cases}
\displaystyle
\exp\!\left[-\left(
\frac{d_{ij}}{\sqrt{\sigma_i\sigma_j}+\eps}
\right)^2\right],&j\in\mathcal N_p(i),\\[2mm]
0,&\text{otherwise}.
\end{cases}
\end{equation}
Set
\[
\widehat w_{ij}^{(1)}=\max\{a_{ij},a_{ji}\},
\qquad
\overline w=\frac{1}{|\mathcal E_p|}\sum_{\{r,s\}\in\mathcal E_p}\widehat w_{rs}^{(1)},
\]
and define
\begin{equation}
\label{eq:local_graph}
w_{ij}^{(1)}=
\begin{cases}
\widehat w_{ij}^{(1)}/\overline w,&\{i,j\}\in\mathcal E_p,\\
0,&\text{otherwise},
\end{cases}
\qquad
w_{ii}^{(1)}=0.
\end{equation}
Therefore, $W^{(0)}$ and $W^{(1)}$ are symmetric, have identical off-diagonal support, and differ only in their edge weights.

\subsection{Out-of-fold graph assessment}

Partition the labeled set into $F$ nonempty approximately stratified folds,
\[
\mathcal L=\mathcal L_1\mathbin{\dot\cup}\cdots\mathbin{\dot\cup}\mathcal L_F.
\]
Here, approximate stratification balances the available labeled samples from each class across the folds as evenly as their integer counts permit. It does not require every class to occur in every held-out fold. This distinction is necessary when a class contains fewer than $F$ labeled samples. Each labeled sample is held out exactly once, and each training set retains the remaining labeled samples from that class whenever the class contains at least two labeled observations.
For fold $f$, define
\[
\mathcal T_f=\mathcal L\setminus\mathcal L_f,
\qquad
\mathcal Q_f=\{1,\ldots,n\}\setminus\mathcal T_f.
\]
For $q\in\{0,1\}$, let
\[
L^{(q)}=D^{(q)}-W^{(q)},
\qquad
D^{(q)}=\operatorname{diag}(W^{(q)}\ones).
\]
The harmonic label scores $H_f^{(q)}\in\R^{|\mathcal Q_f|\times c}$ satisfy
\begin{equation}
\label{eq:harmonic}
L_{\mathcal Q_f\mathcal Q_f}^{(q)}H_f^{(q)}
=W_{\mathcal Q_f\mathcal T_f}^{(q)}C_{\mathcal T_f},
\end{equation}
where $C_{\mathcal T_f}$ contains the one-hot labels in $\mathcal T_f$. If MATLAB reports a singular system or the solution contains nonfinite entries, the stabilized system with $10^{-10}I$ is solved. Negative roundoff values are truncated to zero, and each row with positive sum is normalized. A row with zero sum is retained as the zero score vector. Such a vector is not interpreted as a probability distribution; in \eqref{eq:brier} it receives unit squared loss against any one-hot target and therefore penalizes a candidate graph that cannot propagate labeled information to that observation. Denote the resulting score vector for $i\in\mathcal L_f$ by $p_{f,i}^{(q)}$. The foldwise squared prediction loss is
\begin{equation}
\label{eq:brier}
R_f^{(q)}=
\frac{1}{|\mathcal L_f|}
\sum_{i\in\mathcal L_f}
\left\|p_{f,i}^{(q)}-e_{y_i}\right\|_2^2.
\end{equation}
Define
\begin{equation}
\label{eq:advantage}
\delta_f=R_f^{(0)}-R_f^{(1)},
\qquad
\Delta=\frac{1}{F}\sum_{f=1}^{F}\delta_f,
\end{equation}
and
\begin{equation}
\label{eq:standard_error}
s=\left[
\frac{1}{F(F-1)}
\sum_{f=1}^{F}(\delta_f-\Delta)^2
\right]^{1/2}.
\end{equation}
Thus, \(\Delta>0\) favors the locally weighted graph, whereas \(s\)
summarizes the fold-to-fold variability of this advantage. It is used
as a variability measure rather than a confidence bound. All held-out predictions use labels from $\mathcal T_f$ only; no ground-truth label from $\mathcal U$ is used.

\subsection{Variability-adjusted conservative calibration}

To make the role of foldwise variability explicit,  consider the shrinkage family:
\begin{equation}
\label{eq:theta_kappa}
\theta_{\kappa}=
\begin{cases}
\displaystyle\frac{(\Delta-\kappa s)_{+}}{\Delta},&\Delta>0,\\[1mm]
0,&\Delta\le0,
\end{cases}
\quad \kappa\ge0,
\end{equation}
where $(a)_{+}:=\max\{a,0\}$.
Here, \(\kappa\) controls how strongly the foldwise variability is
discounted. When $\kappa=0$, $\theta_{\kappa}=1$ for every $\Delta>0$, which is the hard graph-selection rule. For $\kappa>0$, a positive advantage is retained only to the extent that it exceeds $\kappa s$, and exact fallback occurs whenever $\Delta\le\kappa s$. Both formal CGC instantiations use $F=5$ and $\kappa=1$, so
\begin{equation}
\label{eq:theta}
\theta:=\theta_1=
\begin{cases}
\displaystyle\frac{(\Delta-s)_{+}}{\Delta},&\Delta>0,\\[2mm]
0,&\Delta\le0.
\end{cases}
\end{equation}
Thus, \((\Delta-\kappa s)_{+}\) is the variability-adjusted retained advantage, and $\theta_\kappa$ is the fraction of the observed advantage that remains after the stability discount. Section~\ref{sec:experiments} examines $\kappa\in\{0,0.5,1,1.5\}$ and $F\in\{3,5,10\}$.

\begin{proposition}[Basic properties of the calibration coefficient]
\label{prop:theta_properties}
Let $s\ge0$, $\kappa\ge0$, and let $\theta_\kappa$ be defined by \eqref{eq:theta_kappa}. Then $0\le\theta_\kappa\le1$. For fixed $\Delta>0$, $\theta_\kappa$ is nonincreasing in $s$ and $\kappa$; for fixed $s$ and $\kappa$, it is nondecreasing in $\Delta>0$. Moreover, for fixed $\Delta>0$,
\[
|\theta_\kappa(s_1)-\theta_\kappa(s_2)|
\le \frac{\kappa}{\Delta}|s_1-s_2|.
\]
In particular, $\theta_\kappa=0$ whenever $\Delta\le\kappa s$.
\end{proposition}

\begin{proof}
For $\Delta>0$,
\[
\theta_\kappa=\max\!\left\{1-\frac{\kappa s}{\Delta},0\right\}.
\]
The range and monotonicity statements follow directly. Since the positive-part map is $1$-Lipschitz, the stated bound follows immediately. For $\Delta\le0$, $\theta_\kappa=0$ by definition.
\end{proof}

The calibrated graph is
\begin{equation}
\label{eq:calibrated_graph}
W_{\theta}=(1-\theta)W^{(0)}+\theta W^{(1)}.
\end{equation}
Because $0\le\theta\le1$, this is a convex combination of two symmetric nonnegative graphs with identical support. Let
\[
D_{\theta}=\operatorname{diag}(W_{\theta}\ones),
\qquad
L_{\theta}=D_{\theta}-W_{\theta}.
\]

\begin{proposition}[Controlled structural deviation]
\label{prop:structural_deviation}
For every $\theta\in[0,1]$,
\[
\begin{aligned}
W_\theta-W^{(0)}
&=\theta\bigl(W^{(1)}-W^{(0)}\bigr),\\
L_\theta-L^{(0)}
&=\theta\bigl(L^{(1)}-L^{(0)}\bigr).
\end{aligned}
\]
Consequently, for any absolutely homogeneous matrix norm,
\[
\begin{aligned}
\|W_\theta-W^{(0)}\|
&=\theta\|W^{(1)}-W^{(0)}\|,\\
\|L_\theta-L^{(0)}\|
&=\theta\|L^{(1)}-L^{(0)}\|.
\end{aligned}
\]
\end{proposition}

\begin{proof}
The first identity follows from \eqref{eq:calibrated_graph}. Because the degree operator is linear,
\[
D_\theta=(1-\theta)D^{(0)}+\theta D^{(1)},
\]
and therefore $L_\theta=(1-\theta)L^{(0)}+\theta L^{(1)}$. The norm equalities follow from $\theta\ge0$ and absolute homogeneity.
\end{proof}

\subsection{Primary instantiation in one-hot-constrained NMF}
\label{subsec:goc_instantiation}

We now instantiate CGC in GOCNMF. The parent model solves
\begin{eqnarray}
\min_{U,V_{\mathcal U:}\ge0}
&\|X-UV^{T}\|_{F}^{2}
+\alpha\Tr(V^{T}L^{(0)}V),\nonumber\\
&V_{\mathcal L:}=C_{\mathcal L}.
\end{eqnarray}
where $U\in\R_{+}^{m\times c}$, $V\in\R_{+}^{n\times c}$, and $\alpha>0$. The unlabeled samples are assigned directly by
\begin{equation}
\label{eq:direct_assignment}
\widehat y_i=\argmaxop_{1\le k\le c}v_{ik},
\qquad i\in\mathcal U.
\end{equation}
For this model, \eqref{eq:generic_graph_objective} has $Z=(U,V)$, $H(Z)=V$, $\gamma=\alpha$, and $\Phi(U,V)=\|X-UV^T\|_F^2$ subject to the one-hot labeled rows. Replacing $L^{(0)}$ by $L_\theta$ gives CGC-GOCNMF:
\begin{eqnarray}
\min_{U,V_{\mathcal U:}\ge0}
J_{\theta}(U,V)
&=&\|X-UV^{T}\|_{F}^{2}
+\alpha\Tr(V^{T}L_{\theta}V),\nonumber\\
V_{i:}&=&e_{y_i}^{T},\quad i\in\mathcal L.\label{eq:cgc_objective}
\end{eqnarray}

\begin{proposition}[Controlled downstream objective perturbation]
\label{prop:objective_deviation}
Let $\mathcal J_L$ be any objective of the form \eqref{eq:generic_graph_objective}. For every feasible $Z$,
\[
\mathcal J_{L_\theta}(Z)-\mathcal J_{L^{(0)}}(Z)
=
\gamma\theta\Tr\!\left[
H(Z)^T\bigl(L^{(1)}-L^{(0)}\bigr)H(Z)
\right],
\]
and therefore
\[
\left|
\mathcal J_{L_\theta}(Z)-\mathcal J_{L^{(0)}}(Z)
\right|
\le
\gamma\theta
\|L^{(1)}-L^{(0)}\|_2
\|H(Z)\|_F^2.
\]
\end{proposition}

\begin{proof}
The graph-independent term $\Phi(Z)$ cancels. Proposition~\ref{prop:structural_deviation} gives $L_\theta-L^{(0)}=\theta(L^{(1)}-L^{(0)})$, which yields the identity. Since $L^{(1)}-L^{(0)}$ is symmetric,
\[
\left|
\Tr\!\left[
H^T(L^{(1)}-L^{(0)})H
\right]
\right|
\le
\|L^{(1)}-L^{(0)}\|_2\|H\|_F^2.
\]
This proof is complete.
\end{proof}

The calibration coefficient is fixed before this optimization and is not updated with $U$ or $V$. Thus, CGC changes only the Laplacian supplied to the parent objective: the one-hot constraints, nonnegativity conditions, factorization variables, and direct assignment rule remain unchanged. The choice $\kappa=1$ belongs to the calibration stage rather than to the NMF regularization parameters.

For fixed $W_{\theta}$, the multiplicative updates are
\begin{equation}
\label{eq:update_u}
U\leftarrow U\had\frac{XV}{U(V^{T}V)}
\end{equation}
and
\begin{equation}
\label{eq:update_v}
V_{\mathcal U:}\leftarrow V_{\mathcal U:}\had
\frac{(X^{T}U+\alpha W_{\theta}V)_{\mathcal U:}}
{(V(U^{T}U)+\alpha D_{\theta}V)_{\mathcal U:}},
~
V_{\mathcal L:}=C_{\mathcal L}.
\end{equation}
Fractions and products are entrywise where appropriate. The analysis in Section~\ref{sec:analysis} concerns the exact
multiplicative updates under the strict-positivity condition stated in
Theorem~\ref{thm:monotonicity}. In the implementation, denominators are lower bounded by
\(\eps\) for numerical protection.

\medskip
\noindent\rule{\linewidth}{0.4pt}
\noindent\textbf{Algorithm 1.} CGC instantiated in GOCNMF (CGC-GOCNMF)
\par\noindent\rule{\linewidth}{0.4pt}

\noindent\textbf{Input:} $X$, labeled set $\mathcal L$, labels $\{y_i:i\in\mathcal L\}$, number of classes $c$, neighbors $p$, graph coefficient $\alpha$, folds $F=5$, maximum iterations $T$.

\noindent\textbf{Output:} $U$, $V$, calibration coefficient $\theta$, and assignments $\widehat y$.

\begin{enumerate}[leftmargin=2em]
\item Normalize each column of $X$ to unit Euclidean norm.
\item Construct $W^{(0)}$ from \eqref{eq:binary_graph} and $W^{(1)}$ from \eqref{eq:local_scale}-\eqref{eq:local_graph}.
\item Divide $\mathcal L$ into $F=5$ nonempty approximately stratified folds, balancing each class as evenly as its labeled count permits.
\item For $q=0,1$ and $f=1,\ldots,F$, hide $\mathcal L_f$, solve the harmonic propagation problem, and compute $R_f^{(q)}$.
\item Compute $\delta_f$, $\Delta$, and $s$, set $\kappa=1$, and obtain $\theta$ and $W_{\theta}$ from \eqref{eq:theta}-\eqref{eq:calibrated_graph}.
\item Initialize positive $U$ and $V_{\mathcal U:}$, and fix $V_{\mathcal L:}=C_{\mathcal L}$.
\item Repeat \eqref{eq:update_u} and \eqref{eq:update_v} for $T$ iterations.
\item Return $\widehat y_i=\argmaxop_{1\le k\le c}v_{ik}$ for $i\in\mathcal U$.
\end{enumerate}

\noindent\rule{\linewidth}{0.4pt}

\subsection{Second downstream instantiation in GNMFLD}
\label{subsec:gnmfld_instantiation}

To test whether CGC depends on the hard one-hot rows of GOCNMF, we also instantiate it in GNMFLD \cite{Xing2021GNMFLD}. Let $Y\in\R_{+}^{n\times c}$ contain the available labels, with $Y_{i:}=e_{y_i}^{T}$ for $i\in\mathcal L$ and $Y_{i:}=0$ otherwise, and let $\mathcal P_{\mathcal L}$ retain the rows indexed by $\mathcal L$ and set all remaining rows to zero. The original GNMFLD label-discrimination term is $\|\mathcal P_{\mathcal L}(V)-\mathcal P_{\mathcal L}(Y)\|_F^2$; here $\mathcal P_{\mathcal L}(Y)=Y$ by construction. The parent objective is therefore
\begin{eqnarray}
&&\min_{U,V\ge0}
\|X-UV^{T}\|_{F}^{2}
+\alpha\|\mathcal P_{\mathcal L}(V)-\mathcal P_{\mathcal L}(Y)\|_{F}^{2}\nonumber\\
&&~~~~~~~~+\beta\Tr(V^{T}L^{(0)}V).\label{eq:gnmfld_parent}
\end{eqnarray}
where $\alpha,\beta>0$. Thus $\alpha$ weights label discrimination and $\beta$ weights graph regularization. Unlike GOCNMF, the labeled rows of $V$ are not fixed; they are encouraged toward their class indicators by the discrimination term.

CGC-GNMFLD replaces only $L^{(0)}$ by $L_\theta$:
\begin{eqnarray}
\min_{U,V\ge0}
J_{\theta}^{\mathrm{LD}}(U,V)
&=&\|X-UV^{T}\|_{F}^{2}\nonumber\\
&&+\alpha\|\mathcal P_{\mathcal L}(V)-\mathcal P_{\mathcal L}(Y)\|_{F}^{2}\nonumber\\
&&+\beta\Tr(V^{T}L_{\theta}V).
\end{eqnarray}
Thus \eqref{eq:generic_graph_objective} applies with $H(Z)=V$ and $\gamma=\beta$. Proposition~\ref{prop:objective_deviation} gives
\[
\left|
J_{\theta}^{\mathrm{LD}}(U,V)-J_{0}^{\mathrm{LD}}(U,V)
\right|
\le
\beta\theta
\|L^{(1)}-L^{(0)}\|_2
\|V\|_F^2.
\]
For a fixed calibrated graph, the parent multiplicative updates are retained:
\[
U\leftarrow U\had\frac{XV}{U(V^{T}V)},
\]
\[
V\leftarrow V\had
\frac{X^{T}U+\alpha Y+\beta W_{\theta}V}
{V(U^{T}U)+\alpha\mathcal P_{\mathcal L}(V)+\beta D_{\theta}V}.
\]
where $\mathcal P_{\mathcal L}(Y)=Y$. The unlabeled assignments remain $\widehat y_i=\argmaxop_k v_{ik}$. No additional coefficient is introduced: the transfer uses the same $F=5$, $\kappa=1$ calibration rule and changes only the graph supplied to the GNMFLD solver.

\section{Structural properties and primary-instantiation convergence}
\label{sec:analysis}

\subsection{Validity and support preservation}

\begin{proposition}[Support preservation]
The two candidate graphs satisfy
\[
w_{ij}^{(0)}>0\quad\Longleftrightarrow\quad w_{ij}^{(1)}>0.
\]
Consequently, $W_{\theta}$ has the same off-diagonal support for every $\theta\in[0,1]$.
\end{proposition}

\begin{proof}
Every directed nearest-neighbor edge receives a strictly positive exponential weight in \eqref{eq:directed_local_graph}. Max-symmetrization therefore gives a positive entry exactly on the union of the directed neighbor relations, which is $\mathcal E_p$. Rescaling by the positive mean $\overline w$ does not change the support. A convex combination of two positive entries remains positive, while two zero entries remain zero. This completes the proof.
\end{proof}

\begin{theorem}[Graph validity]
For every $\theta\in[0,1]$, $W_{\theta}$ is symmetric and nonnegative, and $L_{\theta}$ is positive semidefinite.
\end{theorem}

\begin{proof}
Both candidate graphs are symmetric and nonnegative, so their convex combination has the same properties. For any $z\in\R^n$,
\[
z^{T}L_{\theta}z=\frac12\sum_{i=1}^{n}\sum_{j=1}^{n}w_{\theta,ij}(z_i-z_j)^2\ge0.
\]
Thus $L_{\theta}\succeq0$. The proof is complete.
\end{proof}

\subsection{Exact fallback}

\begin{theorem}[Exact baseline recovery]
\label{thm:exact_recovery}
If $\Delta\le s$, then $\theta=0$, $W_{\theta}=W^{(0)}$, and $L_{\theta}=L^{(0)}$. Consequently, every downstream model of the form \eqref{eq:generic_graph_objective} exactly recovers its baseline graph-regularized objective.
\end{theorem}

\begin{proof}
If $\Delta\le0$, the second branch of \eqref{eq:theta} gives $\theta=0$. If $0<\Delta\le s$, then $1-s/\Delta\le0$, so the first branch also gives $\theta=0$. Equation~\eqref{eq:calibrated_graph} yields $W_\theta=W^{(0)}$, and linearity of the degree operator gives $L_\theta=L^{(0)}$. Substitution into \eqref{eq:generic_graph_objective} proves exact baseline recovery.
\end{proof}

\begin{corollary}[Exact recovery of the instantiated parents]
Under the condition of Theorem~\ref{thm:exact_recovery}, CGC-GOCNMF and CGC-GNMFLD exactly recover their respective parent objectives. With the same labeled split, initialization, solver, and iteration budget, the corresponding multiplicative updates, learned representations, and direct assignments also coincide.
\end{corollary}

\begin{remark}
Exact recovery is a property of the calibration stage itself and is independent of whether the parent uses hard one-hot rows or a label-discrimination penalty.
\end{remark}

\subsection{Monotonicity of the primary GOCNMF instantiation}

\begin{theorem}[Monotonic objective values]\label{thm:monotonicity}
Assume that \(W_{\theta}\) is fixed before factorization and that,
throughout the iterations, all free entries of \(U\) and
\(V_{\mathcal U:}\), together with the denominators in
\eqref{eq:update_u}-\eqref{eq:update_v}, remain strictly positive.
With \(V_{\mathcal L:}\) kept fixed, one complete alternating iteration
satisfies
\[
J_{\theta}\bigl(U^{(t+1)},V^{(t+1)}\bigr)
\le
J_{\theta}\bigl(U^{(t)},V^{(t)}\bigr).
\]
Consequently, the sequence of objective values converges.
\end{theorem}

\begin{proof}
We first use the following elementary inequality. If $M$ is symmetric and entrywise nonnegative, $z>0$, and $h$ is any real vector, then
\begin{equation}
\label{eq:majorization_inequality}
h^{T}Mh\le\sum_k\frac{(Mz)_k}{z_k}h_k^2.
\end{equation}
Indeed,
\begin{eqnarray*}
&&\sum_k\frac{(Mz)_k}{z_k}h_k^2-h^{T}Mh\\
&=&\frac12\sum_{k,\ell}m_{k\ell}
\left(\sqrt{\frac{z_{\ell}}{z_k}}h_k-\sqrt{\frac{z_k}{z_{\ell}}}h_{\ell}\right)^2\ge0.
\end{eqnarray*}

We now consider the two block updates separately. For fixed $V$, let
\[
C=V^{T}V,\qquad E=XV,
\]
and denote by $H=U^{+}-U$ the change produced by \eqref{eq:update_u}. Since the graph term is independent of $U$,
\[
J_{\theta}(U^{+},V)-J_{\theta}(U,V)
=2\langle UC-E,H\rangle+\Tr(HCH^{T}).
\]
Because $C$ is symmetric and entrywise nonnegative, \eqref{eq:majorization_inequality}, applied to every row of $U$, gives
\[
\Tr(HCH^{T})\le\sum_{r,k}\frac{(UC)_{rk}}{u_{rk}}h_{rk}^2.
\]
Moreover, update \eqref{eq:update_u} implies
\[
h_{rk}=-u_{rk}\frac{(UC-E)_{rk}}{(UC)_{rk}},
\]
and
\[
2(UC-E)_{rk}h_{rk}=-2\frac{(UC)_{rk}}{u_{rk}}h_{rk}^2.
\]
Therefore,
\[
J_{\theta}(U^{+},V)-J_{\theta}(U,V)
\le-\sum_{r,k}\frac{(UC)_{rk}}{u_{rk}}h_{rk}^2\le0.
\]

Next fix $U=U^{+}$. Set
\[
A=U^{T}U,\qquad B=X^{T}U,
\]
and define
\[
P=VA+\alpha D_{\theta}V,
\qquad
Q=B+\alpha W_{\theta}V.
\]
Let $G=V^{+}-V$ and $G_{\mathcal L:}=0$, because the labeled rows remain fixed. Expanding the quadratic objective gives
\begin{eqnarray}
\label{eq:v_block_expansion}
J_{\theta}(U,V^{+})-J_{\theta}(U,V)
&=&2\langle P-Q,G\rangle+\Tr(GAG^{T})\nonumber\\
&&+\alpha\Tr(G^{T}L_{\theta}G).
\end{eqnarray}
Since $A$ is symmetric and entrywise nonnegative, \eqref{eq:majorization_inequality} yields
\[
\Tr(GAG^{T})\le\sum_{i\in\mathcal U}\sum_k\frac{(VA)_{ik}}{v_{ik}}g_{ik}^2.
\]
Furthermore, writing $g_i$ for the $i$-th row of $G$, graph symmetry and nonnegativity give
\[
\begin{aligned}
\Tr(G^{T}L_{\theta}G)
&=\frac12\sum_{i,j}w_{\theta,ij}\|g_i-g_j\|_2^2\\
&\le \sum_{i,j}w_{\theta,ij}
\bigl(\|g_i\|_2^2+\|g_j\|_2^2\bigr)\\
&=2\sum_i d_{\theta,i}\|g_i\|_2^2\\
&=2\sum_{i\in\mathcal U}\sum_k
\frac{(D_{\theta}V)_{ik}}{v_{ik}}g_{ik}^2.
\end{aligned}
\]
By \eqref{eq:update_v}, we have
\[
g_{ik}=-v_{ik}\frac{(P-Q)_{ik}}{P_{ik}},
\qquad i\in\mathcal U,
\]
and
\[
2(P-Q)_{ik}g_{ik}=-2\frac{P_{ik}}{v_{ik}}g_{ik}^2.
\]
Substituting these estimates into \eqref{eq:v_block_expansion}, we obtain that
\begin{eqnarray*}
&&J_{\theta}(U,V^{+})-J_{\theta}(U,V)\\
&\le& -2\sum_{i\in\mathcal U}\sum_k
\frac{P_{ik}}{v_{ik}}g_{ik}^2+\sum_{i\in\mathcal U}\sum_k
\frac{(VA)_{ik}}{v_{ik}}g_{ik}^2\\
&&+2\alpha\sum_{i\in\mathcal U}\sum_k
\frac{(D_{\theta}V)_{ik}}{v_{ik}}g_{ik}^2\\
&=&-\sum_{i\in\mathcal U}\sum_k
\frac{(VA)_{ik}}{v_{ik}}g_{ik}^2\le0.
\end{eqnarray*}
Thus both block updates are nonincreasing, and
\[
J_{\theta}\bigl(U^{(t+1)},V^{(t+1)}\bigr)
\le J_{\theta}\bigl(U^{(t+1)},V^{(t)}\bigr)
\le J_{\theta}\bigl(U^{(t)},V^{(t)}\bigr).
\]
Finally, graph validity gives $\Tr(V^{T}L_{\theta}V)\ge0$, so $J_{\theta}(U,V)\ge0$. Hence, the nonincreasing objective sequence is bounded below and therefore converges.
This completes the proof.
\end{proof}

\begin{remark}
Theorem~\ref{thm:monotonicity} establishes convergence of the objective values for the exact updates under the stated strict-positivity condition; finite-precision behavior with denominator protection is examined experimentally.
\end{remark}

\begin{proposition}[Stationarity of positive fixed points]
\label{prop:fixed_stationarity}
Let $(U^\star,V^\star)$ be a fixed point of the exact multiplicative updates with all free entries strictly positive and $V^\star_{\mathcal L:}=C_{\mathcal L}$. Then the gradient of $J_\theta$ vanishes with respect to $U$ and the free block $V_{\mathcal U:}$. Hence $(U^\star,V^\star)$ is first-order stationary for the reduced problem with the labeled rows fixed.
\end{proposition}

\begin{proof}
At a strictly positive fixed point, each multiplicative ratio in \eqref{eq:update_u} and \eqref{eq:update_v} equals one. Therefore
\[
U^\star((V^\star)^TV^\star)=XV^\star,
\]
and, on the free rows,
\[
\bigl(V^\star((U^\star)^TU^\star)+\alpha D_\theta V^\star\bigr)_{\mathcal U:}
=
\bigl(X^TU^\star+\alpha W_\theta V^\star\bigr)_{\mathcal U:}.
\]
These are precisely the zero-gradient conditions for $U$ and $V_{\mathcal U:}$ in \eqref{eq:cgc_objective}.
\end{proof}

\subsection{Computational complexity}

Let $|\mathcal E_p|=O(np)$ denote the number of edges in the common $p$-nearest-neighbor support. Constructing the two candidate graphs requires $O(mn^2)$ time and $O(np)$ sparse storage.

For each graph, the calibration stage performs one harmonic propagation solve in each of the $F$ folds. If $\mathcal S(n,|\mathcal E_p|,c)$ denotes the cost of one sparse solve with $c$ right-hand sides, the calibration cost is
\[
O\!\left(2F\,\mathcal S(n,|\mathcal E_p|,c)\right).
\]
The computation of $\Delta$, $s$, and $\theta_{\kappa}$ costs only $O(F)$. In the proposed method, $F=5$ and $\kappa=1$ are fixed; the alternative values used in the sensitivity study are not part of a single execution.

After calibration, $W_{\theta}$ remains fixed. Each multiplicative iteration costs
\[
O(mnc+npc+nc^2),
\]
and $T$ iterations require
\[
O\!\left(T(mnc+npc+nc^2)\right).
\]
Hence, the total cost is
\[
\begin{aligned}
O\!\bigl(&mn^2+2F\,\mathcal S(n,|\mathcal E_p|,c)
+T(mnc+npc+nc^2)\bigr).
\end{aligned}
\]
CGC-GOCNMF therefore has the same factorization complexity as GOCNMF, with additional computation confined to the pre-factorization calibration stage. The same separation holds for CGC-GNMFLD: calibration adds the common pre-optimization cost, while the downstream factorization retains the parent GNMFLD iteration structure.


%%%%%%%%%%%%%%%%%%%%%%%%%%%%%%%%%%%%%%%%%%%%%%%%%%%%%%%%%%%%%%%%%%%%%%%%%%%%%%%%%%
\section{Experimental evaluation and analysis}
\label{sec:experiments}

\subsection{Datasets and experimental protocol}

We evaluated GNMFLD, ERDNMF, GOCNMF, and CGC-GOCNMF on six image datasets: PIE, YaleB, COIL20, COIL100, Optdigits, and MNIST. Table~\ref{tab:datasets} summarizes the dataset statistics and the graph parameters used by GOCNMF and CGC-GOCNMF.

\begin{table*}[!t]
\centering
\caption{Dataset characteristics and fixed graph settings used in the experiments.}
\label{tab:datasets}
\small
\begin{tabular*}{\textwidth}{@{\extracolsep{\fill}}lrrrrrr@{}}
\toprule
Dataset & Samples $n$ & Features $m$ & Classes $c$ & Labels/trial & Neighbors $p$ & $\alpha$ \\
\midrule
PIE & 2856 & 1024 & 68 & 272 & 3 & 1000 \\
YaleB & 2414 & 1024 & 38 & 227 & 2 & 1000 \\
COIL20 & 1440 & 1024 & 20 & 140 & 3 & 10 \\
COIL100 & 7200 & 1024 & 100 & 700 & 3 & 10 \\
Optdigits & 5620 & 64 & 10 & 557 & 4 & 10 \\
MNIST & 6996 & 784 & 10 & 695 & 4 & 10 \\
\bottomrule
\end{tabular*}
\vspace{1.5pt}

\parbox{\textwidth}{\footnotesize\emph{Notes.} Labels/trial is the total number of labeled samples in each trial. The neighborhood size $p$ and graph-regularization parameter $\alpha$ were fixed before the paired comparisons and were shared by GOCNMF and CGC-GOCNMF.}
\end{table*}

Each sample was normalized to unit Euclidean norm. In each trial, $\max\{2,\lfloor0.1n_k\rfloor\}$ samples from class $k$ were labeled and the remainder were evaluated. We used 20 paired trials with seeds 20260617--20260636; labeled subsets and initial factors were matched within each controlled comparison. The unified comparison and ablation used separately recorded initialization streams, so their CGC--GOC differences need not coincide numerically although each table is internally paired. Approximate five-fold stratification was audited over all 600 dataset--seed--fold combinations, and every training fold retained every class.

GNMFLD and ERDNMF used 200 iterations, whereas GOCNMF and CGC-GOCNMF used 50. The GOCNMF parent comparison keeps the solver, initialization, and budget fixed; Section~\ref{subsec:second_parent_transfer} applies the same controlled design to GNMFLD, changing only $W^{(0)}$ to $W_\theta$. Both CGC instantiations use $F=5$, $\kappa=1$, $\eps=10^{-12}$ for multiplicative denominators, and the stated $10^{-10}I$ harmonic stabilization when required. No unlabeled ground-truth labels were used for parameter selection or calibration.

Assignments were obtained directly by row-wise $\argmaxop$. ACC and NMI are reported as mean $\pm$ sample standard deviation over 20 trials. Two-sided paired Wilcoxon tests \cite{Wilcoxon1945} with baseline-wise Holm correction were interpreted together with paired effect sizes and win/tie/loss counts. Comparisons with GNMFLD and ERDNMF in the unified table provide method-level context; the two parent comparisons isolate the calibration step.

\subsection{Unified comparison of four direct clustering methods}

Table~\ref{tab:four_methods} compares GNMFLD, ERDNMF, GOCNMF, and CGC-GOCNMF under the same paired protocol. All methods produce cluster assignments directly from their learned representations, without $K$-means or an external classifier. ERDNMF parameters were fixed before the formal trials using public settings or labeled-only pilot validation.

\begin{widetable}
\centering
\caption{Clustering performance of four direct semi-supervised NMF methods on six datasets.}
\label{tab:four_methods}
\small
\setlength{\tabcolsep}{33.2pt}
\begin{tabular}{llcc}
\toprule
Dataset & Method & ACC (\%) & NMI (\%) \\
\midrule
\multirow{4}{*}{PIE}
& GNMFLD & $93.27\pm0.77$ & $94.41\pm0.39$ \\
& ERDNMF & $77.22\pm1.41$ & $84.95\pm0.98$ \\
& GOCNMF & $\mathbf{93.93\pm1.29}$ & $\mathbf{95.50\pm0.77}$ \\
& CGC-GOCNMF & $\mathbf{93.93\pm1.29}^{\ddagger\S}$ & $\mathbf{95.50\pm0.77}^{\ddagger\S}$ \\
\midrule
\multirow{4}{*}{YaleB}
& GNMFLD & $61.85\pm1.30$ & $62.62\pm1.04$ \\
& ERDNMF & $24.63\pm1.61$ & $38.51\pm0.91$ \\
& GOCNMF & $\mathbf{73.28\pm1.53}$ & $\mathbf{73.43\pm1.07}$ \\
& CGC-GOCNMF & $73.28\pm1.53^{\ddagger\S}$ & $73.43\pm1.07^{\ddagger\S}$ \\
\midrule
\multirow{4}{*}{COIL20}
& GNMFLD & $88.70\pm0.98$ & $91.36\pm0.68$ \\
& ERDNMF & $68.34\pm1.79$ & $75.79\pm1.67$ \\
& GOCNMF & $91.63\pm1.01$ & $93.32\pm0.73$ \\
& CGC-GOCNMF & $\mathbf{93.66\pm1.13}^{\dagger\ddagger\S}$ & $\mathbf{94.78\pm0.70}^{\dagger\ddagger\S}$ \\
\midrule
\multirow{4}{*}{COIL100}
& GNMFLD & $81.61\pm0.70$ & $89.48\pm0.28$ \\
& ERDNMF & $54.90\pm1.16$ & $72.48\pm0.76$ \\
& GOCNMF & $85.02\pm0.70$ & $91.28\pm0.39$ \\
& CGC-GOCNMF & $\mathbf{87.60\pm0.69}^{\dagger\ddagger\S}$ & $\mathbf{92.99\pm0.39}^{\dagger\ddagger\S}$ \\
\midrule
\multirow{4}{*}{Optdigits}
& GNMFLD & $98.40\pm0.14$ & $95.93\pm0.29$ \\
& ERDNMF & $82.27\pm0.97$ & $71.14\pm0.98$ \\
& GOCNMF & $98.37\pm0.13$ & $95.88\pm0.26$ \\
& CGC-GOCNMF & $\mathbf{98.63\pm0.12}^{\dagger\ddagger\S}$ & $\mathbf{96.52\pm0.26}^{\dagger\ddagger\S}$ \\
\midrule
\multirow{4}{*}{MNIST}
& GNMFLD & $93.43\pm0.28$ & $85.53\pm0.44$ \\
& ERDNMF & $56.14\pm1.74$ & $45.81\pm0.90$ \\
& GOCNMF & $93.55\pm0.27$ & $85.73\pm0.46$ \\
& CGC-GOCNMF & $\mathbf{94.29\pm0.23}^{\dagger\ddagger\S}$ & $\mathbf{87.15\pm0.40}^{\dagger\ddagger\S}$ \\
\bottomrule
\end{tabular}

\vspace{1mm}
\parbox{0.98\linewidth}{\footnotesize Values are mean percentages $\pm$ sample standard deviations over 20 paired trials on unlabeled samples. Bold values indicate the highest mean for each dataset and metric. $\dagger$, $\ddagger$, and $\S$ denote significant improvements over GOCNMF, GNMFLD, and ERDNMF, respectively.}
\end{widetable}

CGC-GOCNMF achieves the highest reported mean ACC and NMI on COIL20, COIL100, Optdigits, and MNIST among the four evaluated direct methods. It exactly matches GOCNMF on PIE and shows no meaningful difference on YaleB. Under the reported fixed-iteration protocol, it also outperforms ERDNMF on all six datasets. The paired significance results are reported in Table~\ref{tab:holm}, with the GOCNMF comparison isolating the effect of calibration most directly.

\subsection{Common-stopping robustness and a recent reproducible baseline}
\label{subsec:common_stop}

Fixed iteration budgets can confound model quality with solver speed. We therefore performed an additional common-stopping experiment with GNMFLD, GOCNMF, CGC-GOCNMF, and the recent SNMFWLP method \cite{Jing2026SNMFWLP}. For every method, stopping was declared when the relative Frobenius change of the row-normalized representation was at most $10^{-4}$ for five consecutive iterations, after at least 20 iterations, with a common cap of 1000 iterations. Row normalization makes the criterion invariant to representation scaling. The reconstruction or model objective was monitored separately and was not used to stop the solver.

The public SNMFWLP repository contains the core multiplicative updates but no complete experiment driver. We therefore provide a MATLAB R2019a integration wrapper implementing data normalization, graph construction, stratified labeled splits, direct row-wise assignment, the common stopping rule, checkpointing, and CSV output. Its coefficients were selected from $\{1,10,100\}^2$ by a disjoint labeled-only pilot split with seed 20260717; held-out labeled samples did not enter training, and no unlabeled ground-truth label was used. The selected $(\alpha,\beta)$ values were $(100,100)$ for PIE and $(10,100)$ for the other five datasets. Formal runs used seeds 20260617--20260619. The algebraically equivalent association $U(U^TXV)$ was used instead of $(UU^T)XV$ to avoid a dense feature-by-feature intermediate; this changes storage cost, not the update equation.

\begin{widetable}
\centering
\caption{Common-stopping robustness over three disjoint formal seeds. Values are mean ACC/NMI percentages $\pm$ sample standard deviation on unlabeled samples.}
\label{tab:common_stop}
\small
\setlength{\tabcolsep}{5.0pt}
\begin{tabular}{lcccc}
\toprule
Dataset & GNMFLD & GOCNMF & CGC-GOCNMF & SNMFWLP (2026) \\
\midrule
PIE & $92.87\pm1.03/94.25\pm0.56$ & $95.16\pm0.97/96.24\pm0.48$ & $95.16\pm0.97/96.24\pm0.48$ & $96.07\pm0.66/96.69\pm0.42$ \\
YaleB & $61.01\pm0.99/61.92\pm0.56$ & $72.66\pm1.05/73.33\pm0.57$ & $72.69\pm1.04/73.33\pm0.57$ & $61.52\pm0.94/62.42\pm0.33$ \\
COIL20 & $89.23\pm1.36/91.64\pm0.88$ & $92.51\pm0.51/93.91\pm0.18$ & $94.38\pm1.42/95.50\pm0.85$ & $91.90\pm1.09/93.76\pm0.64$ \\
COIL100 & $81.69\pm0.49/89.45\pm0.14$ & $86.47\pm0.10/92.16\pm0.14$ & $89.42\pm0.44/94.22\pm0.17$ & $86.84\pm0.52/92.83\pm0.29$ \\
Optdigits & $98.49\pm0.13/96.08\pm0.26$ & $98.36\pm0.06/95.82\pm0.10$ & $98.60\pm0.05/96.44\pm0.08$ & $98.70\pm0.17/96.61\pm0.39$ \\
MNIST & $93.53\pm0.52/85.68\pm0.79$ & $93.69\pm0.50/85.96\pm0.84$ & $94.30\pm0.37/87.21\pm0.65$ & $94.16\pm0.57/87.10\pm0.92$ \\
\bottomrule
\end{tabular}
\end{widetable}

All 54 GNMFLD/GOCNMF/CGC-GOCNMF runs satisfied the common stopping rule and the independently monitored objective nonincrease audit. CGC-GOCNMF retains its parent-level fallback on PIE and YaleB and remains better than GOCNMF on COIL20, COIL100, Optdigits, and MNIST. SNMFWLP is competitive on PIE, Optdigits, and MNIST, but is not uniformly superior: CGC-GOCNMF has higher means on YaleB, COIL20, and COIL100. SNMFWLP reached the common stopping threshold in all 18 formal runs, while its monitored unconstrained objective was not monotone on five of the six datasets; these diagnostics and its 1000-iteration cap are reported rather than interpreted as a stationary-point convergence guarantee. Because this robustness experiment has three formal seeds, it supports solver-budget robustness but does not replace the 20-seed primary comparison.

\subsection{Paired statistical analysis}

Two-sided paired Wilcoxon signed-rank tests were applied to the 20 trial-wise differences between CGC-GOCNMF and each baseline. Holm correction was performed separately for each baseline over the 12 comparisons formed by six datasets and two metrics. No $p$-value was reported when all paired differences were exactly zero.

\begin{widetable}
\centering
\caption{Holm-adjusted $p$-values for CGC-GOCNMF versus the three baselines.}
\label{tab:holm}
\small
\setlength{\tabcolsep}{10.7pt}
\begin{tabular}{lcccccc}
\toprule
& \multicolumn{2}{c}{vs. GOCNMF} & \multicolumn{2}{c}{vs. GNMFLD} & \multicolumn{2}{c}{vs. ERDNMF} \\
\cmidrule(lr){2-3}\cmidrule(lr){4-5}\cmidrule(lr){6-7}
Dataset & ACC & NMI & ACC & NMI & ACC & NMI \\
\midrule
PIE & -- & -- & $\mathbf{0.0382}$ & $\mathbf{2.29\times10^{-5}}$ & $\mathbf{2.29\times10^{-5}}$ & $\mathbf{2.29\times10^{-5}}$ \\
YaleB & $0.6346$ & $0.6346$ & $\mathbf{4.39\times10^{-4}}$ & $\mathbf{2.29\times10^{-5}}$ & $\mathbf{2.29\times10^{-5}}$ & $\mathbf{2.29\times10^{-5}}$ \\
COIL20 & $\mathbf{4.41\times10^{-4}}$ & $\mathbf{2.29\times10^{-5}}$ & $\mathbf{2.29\times10^{-5}}$ & $\mathbf{2.29\times10^{-5}}$ & $\mathbf{1.77\times10^{-4}}$ & $\mathbf{2.29\times10^{-5}}$ \\
COIL100 & $\mathbf{1.91\times10^{-5}}$ & $\mathbf{1.91\times10^{-5}}$ & $\mathbf{4.39\times10^{-4}}$ & $\mathbf{2.29\times10^{-5}}$ & $\mathbf{2.29\times10^{-5}}$ & $\mathbf{2.29\times10^{-5}}$ \\
Optdigits & $\mathbf{4.41\times10^{-4}}$ & $\mathbf{1.91\times10^{-5}}$ & $\mathbf{4.39\times10^{-4}}$ & $\mathbf{2.29\times10^{-5}}$ & $\mathbf{1.77\times10^{-4}}$ & $\mathbf{2.29\times10^{-5}}$ \\
MNIST & $\mathbf{4.41\times10^{-4}}$ & $\mathbf{1.91\times10^{-5}}$ & $\mathbf{4.39\times10^{-4}}$ & $\mathbf{2.29\times10^{-5}}$ & $\mathbf{2.29\times10^{-5}}$ & $\mathbf{2.29\times10^{-5}}$ \\
\bottomrule
\end{tabular}%

\vspace{1mm}
\parbox{0.98\linewidth}{\footnotesize Bold values are significant at the $0.05$ family-wise level. A dash denotes an exact tie over all 20 paired trials. Holm correction was applied separately for each baseline over the 12 dataset--metric comparisons.}
\end{widetable}

Relative to GOCNMF, CGC-GOCNMF is significant on COIL20, COIL100, Optdigits, and MNIST for both metrics, exactly tied on PIE, and not significant on YaleB. All comparisons with GNMFLD and ERDNMF are significant under the reported protocol. Table~\ref{tab:paired_effects} complements adjusted $p$-values with 10,000-resample seed-level bootstrap intervals and W/T/L counts. The largest parent gains occur on COIL20 and COIL100; the Optdigits ACC gain is small, PIE is an exact tie, and YaleB shows only one negligible loss among 20 trials.

\begin{widetable}
\centering
\caption{Paired changes of CGC-GOCNMF relative to GOCNMF. Values are mean percentage-point changes with seed-level percentile-bootstrap 95\% confidence intervals over 20 paired trials (10,000 resamples; MATLAB random seed 20261009). W/T/L denotes wins, exact ties, and losses.}
\label{tab:paired_effects}
\small
\setlength{\tabcolsep}{26.2pt}
\begin{tabular}{lccr}
\toprule
Dataset & $\Delta$ACC (95\% CI) & $\Delta$NMI (95\% CI) & W/T/L (ACC; NMI) \\
\midrule
PIE       & $0.000\ [0.000,0.000]$ & $0.000\ [0.000,0.000]$ & $0/20/0;\ 0/20/0$ \\
YaleB     & $-0.002\ [-0.007,0.000]$ & $-0.000\ [-0.001,0.000]$ & $0/19/1;\ 0/19/1$ \\
COIL20    & $2.027\ [1.696,2.362]$ & $1.467\ [1.157,1.776]$ & $20/0/0;\ 19/0/1$ \\
COIL100   & $2.577\ [2.467,2.691]$ & $1.703\ [1.626,1.778]$ & $20/0/0;\ 20/0/0$ \\
Optdigits & $0.260\ [0.231,0.287]$ & $0.632\ [0.572,0.693]$ & $20/0/0;\ 20/0/0$ \\
MNIST     & $0.740\ [0.669,0.813]$ & $1.424\ [1.289,1.555]$ & $20/0/0;\ 20/0/0$ \\
\bottomrule
\end{tabular}
\end{widetable}

For reproducibility and comparison transparency, Table~\ref{tab:parameters} reports the parameters and their provenance for GNMFLD and ERDNMF. All settings were fixed before the formal paired trials, without using ground-truth labels from the unlabeled samples. For YaleB and MNIST, the ERDNMF coefficient was selected by five-fold stratified validation on labeled samples only, using 20 fixed pilot seeds disjoint from the formal trials. Candidate values were $\{1,10,10^2,10^3,10^4\}$; candidates failing the objective nonincrease audit were excluded. No unlabeled ground-truth labels were used.

\begin{widetable}
\centering
\caption{Parameters and provenance of the external baselines used in the unified comparison.}
\label{tab:parameters}
\small
\setlength{\tabcolsep}{16.7pt}
\begin{tabular}{lrrrrll}
\toprule
& \multicolumn{4}{c}{GNMFLD} & \multicolumn{2}{c}{ERDNMF} \\
\cmidrule(lr){2-5}\cmidrule(lr){6-7}
Dataset & $p$ & $\alpha$ & $\beta$ & Parameter provenance & $\alpha$ & Parameter provenance \\
\midrule
PIE & 5 & $10^5$ & 10 & Pre-specified & 400 & Original authors' code \\
YaleB & 5 & $10^5$ & 100 & Pre-specified & 1000 & Labeled-only CV \\
COIL20 & 5 & $10^4$ & 10 & Pre-specified & 800 & Original authors' code \\
COIL100 & 5 & $10^4$ & 10 & Transferred from COIL20 & 1000 & Original authors' code \\
Optdigits & 5 & $10^5$ & 10 & Transferred from MNIST & 600 & Original authors' code \\
MNIST & 5 & $10^5$ & 10 & Pre-specified & 1000 & Labeled-only CV \\
\bottomrule
\end{tabular}%

\vspace{1mm}
\parbox{0.98\linewidth}{\footnotesize ``Pre-specified'' denotes settings fixed before the formal trials; ``Transferred'' denotes settings inherited from a related dataset; ``Original authors' code'' denotes parameters taken from the public ERDNMF implementation; and ``Labeled-only CV'' denotes validation using labeled samples only. Both baselines were run for 200 iterations. The transferred and pilot-selected settings are explicitly identified and are not presented as dataset-wise optimal parameters.}
\end{widetable}
%%%%%%%%%%%%%%%%%%%%%%%%%%%%%%%%%%%%%%%%%%%%%%%%%%%%%%%%%%%%%
\subsection{Transfer to a second graph-regularized parent model}
\label{subsec:second_parent_transfer}

The general formulation in Section~\ref{sec:method} is not tied to the hard one-hot rows of GOCNMF. We therefore transferred the same CGC rule to GNMFLD, whose supervision is imposed through the label-discrimination term in \eqref{eq:gnmfld_parent}. For every dataset and seed, GNMFLD and CGC-GNMFLD used the same labeled subset, $U_0,V_0$, $(p,\alpha,\beta)$ values from Table~\ref{tab:parameters}, and 200 multiplicative iterations; only $W^{(0)}$ was replaced by $W_\theta$. Because unit-normalized data satisfy $\|x_i-x_j\|_2^2=2(1-x_i^Tx_j)$, the Euclidean and cosine-distance $p$-NN orderings coincide, so the $p=5$ parent support is preserved. This transfer was executed as a separate paired experiment family with its own recorded initialization stream; consequently, rounded GNMFLD means can differ slightly from Table~\ref{tab:four_methods}, while every GNMFLD/CGC-GNMFLD pair remains exactly matched.

\begin{widetable}
\centering
\caption{Transfer of CGC to GNMFLD over 20 paired trials. Changes are CGC-GNMFLD minus GNMFLD in percentage points.}
\label{tab:cgc_gnmfld_transfer}
\small
\setlength{\tabcolsep}{9.3pt}
\begin{tabular}{lcccccc}
\toprule
Dataset & $\bar\theta$ & ACC: parent $\rightarrow$ CGC & $\Delta$ACC &
NMI: parent $\rightarrow$ CGC & $\Delta$NMI & W/T/L \\
\midrule
PIE       & 0.855 & $93.29\pm0.78\rightarrow93.94\pm0.80$ & $+0.652$ & $94.41\pm0.39\rightarrow95.10\pm0.43$ & $+0.695$ & $20/0/0$ \\
YaleB     & 0.907 & $61.85\pm1.30\rightarrow64.75\pm1.39$ & $+2.899$ & $62.63\pm1.04\rightarrow65.31\pm1.01$ & $+2.684$ & $20/0/0$ \\
COIL20    & 0.860 & $88.70\pm0.98\rightarrow90.92\pm1.04$ & $+2.215$ & $91.37\pm0.68\rightarrow93.09\pm0.78$ & $+1.725$ & $20/0/0$ \\
COIL100   & 0.944 & $81.60\pm0.71\rightarrow86.17\pm0.59$ & $+4.567$ & $89.48\pm0.27\rightarrow92.53\pm0.31$ & $+3.058$ & $20/0/0$ \\
Optdigits & 0.875 & $98.40\pm0.14\rightarrow98.64\pm0.12$ & $+0.237$ & $95.93\pm0.29\rightarrow96.50\pm0.26$ & $+0.566$ & $20/0/0$ \\
MNIST     & 0.945 & $93.43\pm0.28\rightarrow94.29\pm0.28$ & $+0.856$ & $85.53\pm0.44\rightarrow87.14\pm0.47$ & $+1.605$ & $20/0/0$ \\
\bottomrule
\end{tabular}%

\vspace{1mm}
\parbox{0.98\linewidth}{\footnotesize ACC/NMI entries are mean percentages $\pm$ sample standard deviations. W/T/L is identical for ACC and NMI. All 12 dataset--metric comparisons remain significant after exact two-sided paired Wilcoxon testing with Holm correction ($p_{\mathrm{adj}}=2.29\times10^{-5}$). No transfer trial produced an exact fallback.}
\end{widetable}

GNMFLD keeps all representation rows free and uses a discrimination penalty, unlike the hard one-hot rows of GOCNMF. Yet the unchanged CGC rule improves both metrics in every paired trial; gains range from $0.237$ to $4.567$ points in ACC and from $0.566$ to $3.058$ points in NMI. This supports transfer beyond the GOCNMF supervision mechanism, without implying universal nondegradation: the formal guarantees remain exact parent recovery at $\theta=0$ and controlled objective perturbation for $\theta>0$.

\subsection{Industrial image case study on steel surface defects}

We further examined the calibration behavior on the NEU-CLS steel-surface defect database \cite{Song2013NEU}, which contains 1800 grayscale images from six defect categories. Each image was resized to $32\times32$, vectorized, and normalized to unit Euclidean norm. No pretrained representation, external classifier, or $K$-means post-processing was used. This industrial-data case study tests whether the graph-calibration behavior persists outside the six benchmark datasets.

Label rates of $5\%$, $10\%$, and $20\%$ per class were considered. For each rate, 20 paired trials used the same labeled samples and formal seeds. ACC, NMI, and ARI were evaluated only on the unlabeled samples. The parameters were fixed before the formal trials: $p=5$ and $\alpha=10$ for GOCNMF and CGC-GOCNMF; $p=5$, $\alpha=10^4$, and $\beta=10$ for GNMFLD; and $\alpha=10^3$ for ERDNMF. The ERDNMF coefficient was selected using labeled-only pilot validation, without using unlabeled ground-truth labels.

\begin{widetable}
\centering
\caption{Clustering performance on NEU-CLS at different label rates.}
\label{tab:neu_performance}
\small
\setlength{\tabcolsep}{20.2pt}
\begin{tabular}{clccc}
\toprule
Label rate & Method & ACC (\%) & NMI (\%) & ARI (\%) \\
\midrule
\multirow{4}{*}{$5\%$}
& GNMFLD & $32.34\pm5.24$ & $14.94\pm2.91$ & $7.74\pm4.13$ \\
& ERDNMF & $23.26\pm2.61$ & $7.09\pm1.43$ & $4.54\pm0.91$ \\
& GOCNMF & $32.78\pm5.29$ & $15.26\pm2.81$ & $8.09\pm4.19$ \\
& CGC-GOCNMF & $\mathbf{33.86\pm4.51}^{\dagger\ddagger\S}$ & $\mathbf{16.50\pm2.44}^{\dagger\ddagger\S}$ & $\mathbf{8.72\pm3.97}^{\dagger\ddagger\S}$ \\
\midrule
\multirow{4}{*}{$10\%$}
& GNMFLD & $37.76\pm3.72$ & $17.37\pm2.56$ & $10.01\pm3.31$ \\
& ERDNMF & $24.12\pm2.29$ & $7.61\pm1.97$ & $4.89\pm1.19$ \\
& GOCNMF & $37.51\pm3.75$ & $17.20\pm2.64$ & $9.76\pm3.38$ \\
& CGC-GOCNMF & $\mathbf{40.32\pm3.44}^{\dagger\ddagger\S}$ & $\mathbf{20.56\pm2.52}^{\dagger\ddagger\S}$ & $\mathbf{11.85\pm3.30}^{\dagger\ddagger\S}$ \\
\midrule
\multirow{4}{*}{$20\%$}
& GNMFLD & $42.44\pm2.48$ & $20.18\pm2.12$ & $12.72\pm2.59$ \\
& ERDNMF & $25.59\pm2.13$ & $8.05\pm1.02$ & $5.20\pm0.59$ \\
& GOCNMF & $42.05\pm2.54$ & $19.93\pm2.13$ & $12.31\pm2.67$ \\
& CGC-GOCNMF & $\mathbf{46.90\pm2.38}^{\dagger\ddagger\S}$ & $\mathbf{25.94\pm2.27}^{\dagger\ddagger\S}$ & $\mathbf{16.31\pm2.83}^{\dagger\ddagger\S}$ \\
\bottomrule
\end{tabular}

\vspace{1mm}
\parbox{0.98\linewidth}{\footnotesize Values are mean percentages $\pm$ sample standard deviations over 20 paired trials on unlabeled samples. Bold values indicate the highest mean at each label rate. $\dagger$, $\ddagger$, and $\S$ denote significant improvements over GOCNMF, GNMFLD, and ERDNMF, respectively. Holm correction was applied separately for each baseline over the nine application comparisons.}
\end{widetable}

Table~\ref{tab:neu_performance} shows that CGC-GOCNMF achieves the highest mean ACC, NMI, and ARI among the four evaluated direct methods at all three label rates. Its advantage over GOCNMF increases with the label rate, reaching $4.84$, $6.01$, and $4.00$ percentage points in ACC, NMI, and ARI, respectively, at the $20\%$ label rate. All comparisons with the three baselines remain significant after baseline-wise Holm correction. The main value of this experiment is the consistent calibration behavior across increasing label rates.

\begin{table}[t]
\centering
\caption{Calibration behavior on NEU-CLS over 20 trials.}
\label{tab:neu_calibration}
\small
\begin{tabular}{lrrr}
\toprule
Label rate & Mean $\theta$ & SD of $\theta$ & Exact fallbacks \\
\midrule
$5\%$ & 0.296 & 0.287 & 7/20 \\
$10\%$ & 0.646 & 0.183 & 0/20 \\
$20\%$ & 0.790 & 0.063 & 0/20 \\
\bottomrule
\end{tabular}
\end{table}

As shown in Table~\ref{tab:neu_calibration}, limited labeled evidence produces stronger shrinkage and seven exact fallbacks at the $5\%$ rate. As the label rate increases, $\theta$ rises and becomes more stable, indicating more consistent support for the locally weighted graph.

\begin{figure*}[t]
\centering
\begin{subfigure}[t]{0.48\linewidth}
\centering
\includegraphics[width=\linewidth]{figures/fig-000.png}
\caption{GNMFLD}
\end{subfigure}\hfill
\begin{subfigure}[t]{0.48\linewidth}
\centering
\includegraphics[width=\linewidth]{figures/fig-001.png}
\caption{ERDNMF}
\end{subfigure}

\vspace{2mm}
\begin{subfigure}[t]{0.48\linewidth}
\centering
\includegraphics[width=\linewidth]{figures/fig-002.png}
\caption{GOCNMF}
\end{subfigure}\hfill
\begin{subfigure}[t]{0.48\linewidth}
\centering
\includegraphics[width=\linewidth]{figures/fig-003.png}
\caption{CGC-GOCNMF}
\end{subfigure}
\caption{Aggregated row-normalized confusion matrices on NEU-CLS at the $20\%$ label rate over 20 paired trials: (a) GNMFLD, (b) ERDNMF, (c) GOCNMF, and (d) CGC-GOCNMF. Rows denote true defect categories and columns denote predicted categories. Cr, In, Pa, PS, RS, and Sc denote crazing, inclusion, patches, pitted surface, rolled-in scale, and scratches, respectively. Cell values are percentages, and all panels use the same color scale.}
\label{fig:neu_confusion}
\end{figure*}

Figure~\ref{fig:neu_confusion} presents the aggregated confusion matrices at the $20\%$ label rate. Relative to GOCNMF, CGC-GOCNMF improves the recalls of patches, rolled-in scale, and scratches, although crazing remains difficult when only raw intensity features are used.

\subsubsection{Representation robustness}
To test whether the calibration behavior depends on raw intensity vectors, we repeated the paired GOCNMF/CGC-GOCNMF comparison using a fixed 944-dimensional nonnegative spatial uniform-LBP representation. Each image was resized to $34\times34$, uniform $\mathrm{LBP}_{8,1}^{u2}$ codes were computed on the valid $32\times32$ region, and normalized 59-bin histograms were extracted from a $4\times4$ cell grid. The label rates, 20 formal seeds, labeled subsets, initial factors, $p=5$, $\alpha=10$, 50 factorization iterations, five calibration folds, and direct row-wise assignment were retained.

\begin{widetable}
\centering
\caption{Representation robustness on NEU-CLS using nonnegative spatial uniform-LBP features.}
\label{tab:neu_lbp}
\small
\setlength{\tabcolsep}{11.8pt}
\begin{tabular}{clccccc}
\toprule
Label rate & Method & ACC (\%) & NMI (\%) & ARI (\%) & $\bar{\theta}$ & Fallbacks \\
\midrule
\multirow{2}{*}{$5\%$}
& GOCNMF & $27.42\pm2.37$ & $10.55\pm1.86$ & $3.19\pm2.38$ & -- & -- \\
& CGC-GOCNMF & $27.42\pm2.37$ & $10.55\pm1.86$ & $3.19\pm2.38$ & 0.000 & 20/20 \\
\midrule
\multirow{2}{*}{$10\%$}
& GOCNMF & $32.66\pm1.81$ & $14.51\pm1.13$ & $5.58\pm1.53$ & -- & -- \\
& CGC-GOCNMF & $\mathbf{33.38\pm2.19}^{*}$ & $\mathbf{15.81\pm1.51}^{*}$ & $5.77\pm1.59$ & 0.357 & 8/20 \\
\midrule
\multirow{2}{*}{$20\%$}
& GOCNMF & $38.53\pm1.11$ & $18.74\pm1.40$ & $9.39\pm0.98$ & -- & -- \\
& CGC-GOCNMF & $\mathbf{41.39\pm1.20}^{*}$ & $\mathbf{22.76\pm1.25}^{*}$ & $\mathbf{10.91\pm1.03}^{*}$ & 0.811 & 0/20 \\
\bottomrule
\end{tabular}%

\vspace{1mm}
\parbox{0.98\linewidth}{\footnotesize Values are mean percentages $\pm$ sample standard deviations over 20 paired trials on unlabeled samples. An asterisk denotes a significant improvement over GOCNMF after Holm correction over the nine label-rate--metric comparisons. Fallbacks count trials with $\theta=0$.}
\end{widetable}

At the $5\%$ label rate, CGC-GOCNMF set $\theta=0$ in all 20 trials and exactly reproduced GOCNMF. At the $10\%$ rate, it improved ACC and NMI by $0.72$ and $1.30$ percentage points, respectively; both differences remained significant after Holm correction, whereas the ARI difference was not significant. At the $20\%$ rate, the gains were $2.86$, $4.03$, and $1.52$ percentage points in ACC, NMI, and ARI, respectively, and all three differences were significant. Although the LBP representation did not improve the absolute clustering accuracy over raw pixels, the fallback and gain patterns remained coherent under a qualitatively different nonnegative texture representation.


%%%%%%%%%%%%%%%%%%%%%%%%%%%%%%%%%%%%%%%%%%%%%%%%%%%%%%%%%%%%%%%%%%%%%%%%%%%%%%%%%%%%%%%%%%%%%%%%%%%%%%%%%%%%%%%
\subsection{Six-dataset ablation}

We evaluated BASE, LOCAL, HARD, and SHRINK on all six benchmark datasets. All variants used the same graph support, labeled subsets, initial factors, regularization coefficients, and iteration counts. BASE uses $W^{(0)}$, LOCAL uses $W^{(1)}$, HARD selects $W^{(1)}$ when $\Delta>0$, and SHRINK uses $W_{\theta}$.

\begin{widetable}
\centering
\caption{Six-dataset ablation of local graph weighting and conservative calibration.}
\label{tab:ablation6}
\scriptsize
\setlength{\tabcolsep}{10.2pt}
\begin{tabular}{lccrrcrrcrrc}
\toprule
& & & \multicolumn{3}{c}{LOCAL} & \multicolumn{3}{c}{HARD} & \multicolumn{3}{c}{SHRINK} \\
\cmidrule(lr){4-6}\cmidrule(lr){7-9}\cmidrule(lr){10-12}
Dataset & $\bar\theta$ & Fallbacks & $\Delta$ACC & $\Delta$NMI & Losses & $\Delta$ACC & $\Delta$NMI & Losses & $\Delta$ACC & $\Delta$NMI & Losses \\
\midrule
PIE & 0.000 & 20/20 & -1.558 & -0.761 & 20/20 & 0.000 & 0.000 & 0/0 & 0.000 & 0.000 & 0/0 \\
YaleB & 0.003 & 19/20 & -4.088 & -3.217 & 20/20 & -0.428 & -0.368 & 2/2 & +0.005 & +0.001 & 0/0 \\
COIL20 & 0.770 & 0/20 & +2.958 & +2.171 & 0/0 & +2.958 & +2.171 & 0/0 & +1.881 & +1.420 & 0/0 \\
COIL100 & 0.903 & 0/20 & +2.709 & +1.769 & 0/0 & +2.709 & +1.769 & 0/0 & +2.472 & +1.628 & 0/0 \\
Optdigits & 0.876 & 0/20 & +0.305 & +0.736 & 0/0 & +0.305 & +0.736 & 0/0 & +0.255 & +0.623 & 0/0 \\
MNIST & 0.943 & 0/20 & +0.813 & +1.573 & 0/0 & +0.813 & +1.573 & 0/0 & +0.744 & +1.434 & 0/0 \\
\bottomrule
\end{tabular}%

\vspace{1mm}
\parbox{0.98\linewidth}{\footnotesize Changes are absolute percentage points relative to BASE over 20 paired trials on unlabeled samples. ``Losses'' reports the ACC/NMI loss counts relative to BASE. Fallbacks count trials with $\theta=0$.}
\end{widetable}

Table~\ref{tab:ablation6} reveals two distinct regimes. LOCAL improves all 20 trials on COIL20, COIL100, Optdigits, and MNIST, but degrades both metrics in every PIE and YaleB trial. HARD removes the PIE failure but still selects the local graph in two YaleB trials, producing two ACC and two NMI losses. SHRINK exactly falls back in all PIE trials and in 19 of the 20 YaleB trials, while retaining substantial gains on the other four datasets. No ACC or NMI loss is observed for SHRINK in the 120 paired ablation runs, while the clear-benefit datasets retain most of the local-graph gain.
%%%%%%%%%%%%%%%%%%%%%%%%%%%%%%%%%%%%%%%%%%%%%%%%%%%%%%%%%%%%%%%%%%%%%%%%
\subsection{Sensitivity to the calibration rule}

We generalized the calibration rule to
\[
\theta_{\kappa}=
\begin{cases}
(\Delta-\kappa s)_{+}/\Delta,&\Delta>0,\\
0,&\Delta\le0,
\end{cases}
\]
and considered $\kappa\in\{0,0.5,1,1.5\}$ and $F\in\{3,5,10\}$. PIE, YaleB, and COIL100 were fixed before this experiment to represent fallback, borderline, and clear-benefit regimes, respectively. The labeled subsets and initial factors were paired across all settings.

\begin{widetable}
\centering
\caption{Sensitivity of conservative calibration to the shrinkage strength $\kappa$ and the number of folds $F$.}
\label{tab:sensitivity}
\scriptsize
\setlength{\tabcolsep}{24.2pt}
\begin{tabular}{llrrrrc}
\toprule
Dataset & Setting & $\bar\theta$ & Fallbacks & $\Delta$ACC & $\Delta$NMI & Losses \\
\midrule
\multicolumn{7}{c}{\textit{(a) Sensitivity to $\kappa$ with $F=5$}} \\
\midrule
\multirow{4}{*}{PIE}
& $\kappa=0$ & 0.000 & 20/20 & 0.000 & 0.000 & 0/0 \\
& $\kappa=0.5$ & 0.000 & 20/20 & 0.000 & 0.000 & 0/0 \\
& $\kappa=1$ & 0.000 & 20/20 & 0.000 & 0.000 & 0/0 \\
& $\kappa=1.5$ & 0.000 & 20/20 & 0.000 & 0.000 & 0/0 \\
\addlinespace
\multirow{4}{*}{YaleB}
& $\kappa=0$ & 0.100 & 18/20 & -0.428 & -0.368 & 2/2 \\
& $\kappa=0.5$ & 0.027 & 19/20 & -0.041 & -0.041 & 1/1 \\
& $\kappa=1$ & 0.003 & 19/20 & +0.005 & +0.001 & 0/0 \\
& $\kappa=1.5$ & 0.000 & 20/20 & 0.000 & 0.000 & 0/0 \\
\addlinespace
\multirow{4}{*}{COIL100}
& $\kappa=0$ & 1.000 & 0/20 & +2.709 & +1.769 & 0/0 \\
& $\kappa=0.5$ & 0.952 & 0/20 & +2.685 & +1.774 & 0/0 \\
& $\kappa=1$ & 0.903 & 0/20 & +2.472 & +1.628 & 0/0 \\
& $\kappa=1.5$ & 0.855 & 0/20 & +2.289 & +1.504 & 0/0 \\
\midrule
\multicolumn{7}{c}{\textit{(b) Sensitivity to $F$ with $\kappa=1$}} \\
\midrule
\multirow{3}{*}{PIE}
& $F=3$ & 0.000 & 20/20 & 0.000 & 0.000 & 0/0 \\
& $F=5$ & 0.000 & 20/20 & 0.000 & 0.000 & 0/0 \\
& $F=10$ & 0.000 & 20/20 & 0.000 & 0.000 & 0/0 \\
\addlinespace
\multirow{3}{*}{YaleB}
& $F=3$ & 0.030 & 19/20 & -0.053 & -0.053 & 1/1 \\
& $F=5$ & 0.003 & 19/20 & +0.005 & +0.001 & 0/0 \\
& $F=10$ & 0.001 & 19/20 & +0.002 & -0.001 & 0/1 \\
\addlinespace
\multirow{3}{*}{COIL100}
& $F=3$ & 0.925 & 0/20 & +2.555 & +1.679 & 0/0 \\
& $F=5$ & 0.903 & 0/20 & +2.472 & +1.628 & 0/0 \\
& $F=10$ & 0.908 & 0/20 & +2.525 & +1.665 & 0/0 \\
\bottomrule
\end{tabular}

\vspace{1mm}
\parbox{0.98\linewidth}{\footnotesize Changes are absolute percentage points relative to BASE over 20 paired trials. ``Losses'' reports ACC/NMI loss counts. The proposed setting is $F=5$ and $\kappa=1$.}
\end{widetable}

The qualitative behavior is stable across the tested settings. PIE always returns to BASE, whereas COIL100 retains a clear positive gain for every combination of $F$ and $\kappa$. YaleB exposes the role of conservative shrinkage: with $F=5$, hard selection ($\kappa=0$) produces two losses, while increasing $\kappa$ progressively contracts $\theta$. The proposed choice $\kappa=1$ removes the observed losses while retaining a nonzero coefficient when supported; $\kappa=1.5$ is more conservative and always falls back on YaleB. Varying $F$ has little effect on PIE and COIL100, while the borderline YaleB case is most stable at $F=5$ in these trials. Thus, $F=5$ and $\kappa=1$ provide a balanced conservative setting across the three tested regimes.
%%%%%%%%%%%%%%%%%%%%%%%%%%%%%%%%%%%%%%%%%%%%%%%%%%%%%%%%%%%%%%%%%%%%%%%%%%%%%%%%%%%
\subsection{Controlled graph-reliability stress test}

We next tested whether the calibration reacts to controlled deterioration of the candidate graph. On COIL100, let $W^{(\mathrm{perm})}$ be obtained by randomly permuting the undirected nonzero weights of $W^{(1)}$ over the same edge set. For $\rho\in\{0,0.2,0.4,0.6,0.8,1\}$, define
\[
W^{(\rho)}=(1-\rho)W^{(1)}+\rho W^{(\mathrm{perm})}.
\]
The construction preserves support, symmetry, nonnegativity, and the marginal edge-weight distribution while progressively destroying the association between edge location and local similarity. Performance is reported as paired change relative to BASE.

\begin{widetable}
\centering
\caption{Graph-reliability stress-test results on COIL100.}
\label{tab:stress}
\scriptsize
\setlength{\tabcolsep}{13.2pt}
\begin{tabular}{rrrrrrrrrr}
\toprule
$\rho$ & Edge corr. & $\bar\theta$ & $\bar\Delta$ & \multicolumn{3}{c}{ACC} & \multicolumn{3}{c}{NMI} \\
\cmidrule(lr){5-7}\cmidrule(lr){8-10}
& & & & Candidate & SHRINK & Losses & Candidate & SHRINK & Losses \\
\midrule
0.0 & +1.000 & 0.903 & +0.05736 & +2.709 & +2.472 & 0/0 & +1.769 & +1.628 & 0/0 \\
0.2 & +0.970 & 0.923 & +0.03426 & +2.004 & +1.753 & 0/0 & +1.327 & +1.162 & 0/0 \\
0.4 & +0.832 & 0.915 & +0.02133 & +1.217 & +1.096 & 0/0 & +0.790 & +0.716 & 0/0 \\
0.6 & +0.555 & 0.869 & +0.01185 & +0.508 & +0.455 & 0/1 & +0.329 & +0.302 & 0/0 \\
0.8 & +0.243 & 0.489 & +0.00377 & -0.168 & +0.074 & 14/5 & -0.127 & +0.054 & 14/3 \\
1.0 & -0.000 & 0.018 & -0.00454 & -1.125 & +0.002 & 20/0 & -0.767 & +0.002 & 20/0 \\
\bottomrule
\end{tabular}%

\vspace{1mm}
\parbox{0.98\linewidth}{\footnotesize Results are based on 20 paired trials. Performance changes are absolute percentage points relative to BASE. ``Losses'' reports the candidate/SHRINK loss counts.}
\end{widetable}

\begin{figure}[t]
\centering
\includegraphics[width=0.96\linewidth]{figures/fig-004.png}
\caption{Response of the calibration coefficient to controlled graph corruption on COIL100. Points are means over 20 paired trials, and error bars denote standard errors.}
\label{fig:stress_theta}
\end{figure}

\begin{figure*}[t]
\centering
\begin{subfigure}[t]{0.48\linewidth}
\centering
\includegraphics[width=\linewidth]{figures/fig-005.png}
\caption{ACC}
\end{subfigure}\hfill
\begin{subfigure}[t]{0.48\linewidth}
\centering
\includegraphics[width=\linewidth]{figures/fig-006.png}
\caption{NMI}
\end{subfigure}
\caption{Performance changes relative to BASE under controlled corruption. The uncalibrated candidate deteriorates as $\rho$ increases, whereas SHRINK substantially reduces the losses.}
\label{fig:stress_performance}
\end{figure*}

As $\rho$ increases, the mean edge-weight correlation decreases from 1 to approximately 0, while the mean out-of-fold advantage decreases monotonically from $0.05736$ to $-0.00454$. The uncalibrated candidate changes from a clear improvement to a consistent loss. Although $\bar{\theta}$ is not strictly monotone because it depends on
$s_{\rho}/\Delta_{\rho}$, it decreases to $0.489$ at $\rho=0.8$ and
$0.018$ at $\rho=1$; at $\rho=1$, 19 of the 20 trials yield exact
fallback.

At $\rho=0.8$, the candidate loses ACC and NMI in 14 trials, whereas
SHRINK reduces these counts to 5 and 3. At $\rho=1$, the candidate
loses both metrics in all 20 trials, while SHRINK records no loss and
is practically identical to BASE. These results show that the proposed calibration responds conservatively to severe graph deterioration and progressively returns the downstream model toward the baseline behavior.

%%%%%%%%%%%%%%%%%%%%%%%%%%%%%%%%%%%%%%%%%%%%%%%%%%%%%%%%%%%%%%%%%%%%%%%%
\subsection{Prescribed-protocol training time}

We measured the prescribed-protocol training time of each method over
20 runs. Data loading, common feature normalization, and post-training evaluation were excluded. For the graph-based methods, the measured time included graph construction and factorization. CGC-GOCNMF additionally included the complete five-fold calibration procedure. ERDNMF has no graph or calibration stage, so its reported time consists only of factorization. An untimed warm-up was performed before the formal runs.

\begin{widetable}
\centering
\caption{Prescribed-protocol training times of the four direct methods
in seconds.}
\label{tab:timing}
\scriptsize
\setlength{\tabcolsep}{10.5pt}
\begin{tabular}{lrrrrrrr}
\toprule
Dataset & GNMFLD & ERDNMF & GOCNMF & CGC-GOCNMF & $T_{\mathrm{cal}}$ & Overhead (\%) & Cal. share (\%) \\
\midrule
PIE & $10.719\pm0.490$ & $19.017\pm0.091$ & $4.068\pm0.198$ & $4.504\pm0.426$ & 0.410 & 10.73 & 9.09 \\
YaleB & $5.981\pm0.626$ & $12.049\pm1.581$ & $2.385\pm0.251$ & $2.407\pm0.209$ & 0.050 & 0.95 & 2.10 \\
COIL20 & $3.057\pm0.036$ & $5.236\pm0.073$ & $1.133\pm0.016$ & $1.172\pm0.024$ & 0.024 & 3.41 & 2.07 \\
COIL100 & $49.974\pm7.103$ & $58.937\pm9.199$ & $25.014\pm3.619$ & $28.542\pm3.892$ & 3.704 & 14.11 & 12.98 \\
Optdigits & $8.131\pm0.437$ & $1.875\pm0.028$ & $7.720\pm0.344$ & $7.872\pm0.152$ & 0.176 & 1.98 & 2.24 \\
MNIST & $22.143\pm4.320$ & $24.411\pm2.509$ & $17.492\pm3.066$ & $18.571\pm3.548$ & 0.436 & 6.17 & 2.35 \\
\bottomrule
\end{tabular}%

\vspace{1mm}
\parbox{0.98\linewidth}{\footnotesize $T_{\mathrm{cal}}$ is the CGC-GOCNMF calibration time. ``Overhead'' is the relative increase over GOCNMF, while ``Cal. share'' is the percentage of calibration time in the total CGC-GOCNMF time.}
\end{widetable}

As shown in Table~\ref{tab:timing}, the additional cost of CGC-GOCNMF over GOCNMF ranges from $0.95\%$ to $14.11\%$. Calibration accounts for at most $12.98\%$ of its total running time. Because the methods use different objectives and iteration counts, only
the runtime difference between GOCNMF and CGC-GOCNMF is directly
interpretable; the remaining values are protocol-level references.
%%%%%%%%%%%%%%%%%%%%%%%%%%%%%%%%%%%%%%%%%%%%%%%%%%%%%%%%%%%%%
\subsection{Discussion and limitations}

The experiments reveal complementary calibration regimes rather than a requirement that the candidate graph always be used. In the primary GOCNMF instantiation, stable positive evidence retains substantial local weighting on COIL20, COIL100, Optdigits, and MNIST; weak evidence contracts the coefficient on YaleB; and insufficient evidence produces exact fallback on PIE and in the low-label LBP experiment. The controlled corruption study strengthens this mechanism-level interpretation by showing that deliberate deterioration of edge-location information reduces the estimated candidate advantage and drives the downstream model back toward the baseline.

The GNMFLD transfer addresses a different question: whether this behavior is tied to GOCNMF's hard one-hot supervision. GNMFLD instead uses a label-discrimination penalty, yet the unchanged CGC rule improves both metrics in all 20 paired trials on all six benchmarks. Together, the two instantiations support the intended separation between the calibration stage and the downstream supervision mechanism. Propositions~\ref{prop:theta_properties}--\ref{prop:objective_deviation} explain this separation structurally: $\theta$ controls graph, Laplacian, and downstream graph-regularized objective deviation, while Theorem~\ref{thm:exact_recovery} provides exact parent recovery when the retained evidence vanishes.

Three limitations remain. First, the variability discount is deterministic rather than a calibrated confidence procedure. Second, both downstream instantiations are NMF models and the datasets are predominantly image-based; non-NMF Laplacian-regularized learners and non-image data would test broader transfer. Third, the common-stopping and SNMFWLP robustness study uses three formal seeds, whereas the primary fixed-budget study uses 20; a 20-seed common-stopping replication would permit equally powered inferential comparisons.

%%%%%%%%%%%%%%%%%%%%%%%%%%%%%%%%%%%%%%%%%%%%%%%%%%%%%%%%%%%%%%%%%%%%%%%%%%%%%%
\section{Conclusion}
\label{sec:conclusion}

This paper introduced evidence-guided Conservative Graph Calibration as a pre-optimization mechanism for Laplacian-regularized semi-supervised representation models. Out-of-fold label propagation evaluates a prescribed baseline graph and a same-support locally weighted candidate, while a variability-adjusted retained-advantage coefficient controls how far the final graph departs from the baseline. The analysis shows that this coefficient scales graph and Laplacian deviation, bounds the induced graph-regularized objective perturbation, and yields exact parent recovery when the evidence for modification is insufficient.

The primary CGC-GOCNMF instantiation improves its parent on COIL20, COIL100, Optdigits, and MNIST and retains essentially baseline behavior on PIE and YaleB. The same calibration rule transfers without retuning to GNMFLD, a structurally different parent using label-discrimination supervision: CGC-GNMFLD improves ACC and NMI in all 20 paired trials on each of the six benchmarks, with mean ACC gains of $0.24$--$4.57$ percentage points and NMI gains of $0.57$--$3.06$ points. Industrial-image, representation-robustness, sensitivity, ablation, and controlled-corruption experiments provide complementary evidence about when local weighting is retained, attenuated, or rejected. Taken together, the two parent-model studies support the view that limited labels can be used not only to supervise representation learning but also to regulate structural graph decisions before optimization.



\section*{CRediT authorship contribution statement}

Deshu Sun: Writing-original draft, Methodology, Software, Conceptualization, Formal analysis.  Feng Wang: Writing-review and editing, Funding acquisition, Visualization, Conceptualization. All authors reviewed the manuscript.


\section*{Declaration of competing interest}
The authors declare that they have no known competing financial interests or personal relationships that could have appeared to influence the work reported in this paper.

\section*{Data and code availability}

The MATLAB R2019a code, trial-level outputs, summary files, figures, manuscript source, and the redistributable Optdigits matrix are available at \url{https://github.com/matrix-research-lab-Wf/CGC-NMF}. The repository records exact matrix dimensions and SHA-256 checksums for every dataset used. PIE, YaleB, COIL20, COIL100, MNIST, and NEU-CLS matrices are not redistributed where an explicit compatible redistribution license could not be verified; authoritative access links and local preparation instructions are provided instead.

\section*{Acknowledgments}
This work was supported by Guizhou Provincial Science and Technology Projects (2025039, 2026283), the High-Level Innovative Talent Project of Guizhou Province (GCC2023027), the Natural Science Research Project of Department of Education of Guizhou Province (2024259), and the Research Foundation of Guizhou Minzu University (GZMUZK[2023]YB10).


%%%%%%%%%%%%%%%%%%%%%%%%%%%%%%%%%%%%%%%%%%%%%%%%%%%%%%%%%%%%%%%%%%%

\balance
\footnotesize
\begin{thebibliography}{99}

\bibitem{Lee1999NMF}
Lee, D.~D., Seung, H.~S., 1999.
Learning the parts of objects by non-negative matrix factorization.
\emph{Nature} 401(6755), 788-791.

\bibitem{Lee2000Algorithms}
Lee, D.~D., Seung, H.~S., 2000.
Algorithms for non-negative matrix factorization.
In: \emph{Advances in Neural Information Processing Systems 13}, pp.~556-562.

\bibitem{Cai2011GNMF}
Cai, D., He, X., Han, J., Huang, T.~S., 2011.
Graph regularized nonnegative matrix factorization for data representation.
\emph{IEEE Transactions on Pattern Analysis and Machine Intelligence} 33~(8), 1548-1560.

\bibitem{Liu2012CNMF}
Liu, H., Wu, Z., Li, X., Cai, D., Huang, T.~S., 2012.
Constrained nonnegative matrix factorization for image representation.
\emph{IEEE Transactions on Pattern Analysis and Machine Intelligence} 34~(7), 1299-1311.

\bibitem{Mo2024PNLPSCNMF}
Mo, Y., Li, X., Mei, J., 2024.
Semi-supervised nonnegative matrix factorization with label propagation and constraint propagation.
\emph{Engineering Applications of Artificial Intelligence} 133, 108196.

\bibitem{Xing2021GNMFLD}
Xing, Z., Ma, Y., Yang, X., Nie, F., 2021.
Graph regularized nonnegative matrix factorization with label discrimination for data clustering.
\emph{Neurocomputing} 440, 297-309.

\bibitem{Jing2026SNMFWLP}
Jing, W., Lu, L., Ou, W., 2026.
Semi-supervised non-negative matrix factorization with weighted label propagation for data representation.
\emph{Signal Processing} 241, 110407.

\bibitem{Li2023ABNMTF}
Li, S., Li, W., Lu, H., Li, Y., 2023.
Semi-supervised non-negative matrix tri-factorization with adaptive neighbors and block-diagonal learning.
\emph{Engineering Applications of Artificial Intelligence} 121, 106043.

\bibitem{Ren2025AHSSNMF}
Ren, X., Yang, Y., 2025.
Adaptive hypergraph structure regularized semi-supervised non-negative matrix factorization for image clustering.
\emph{Neurocomputing} 651, 130895.

\bibitem{Yang2025SPNMF}
Yang, X., Zhu, T., Peng, S., Nie, F., Lin, Z., 2025.
Semi-supervised pivotal-aware nonnegative matrix factorization with label and pairwise constraint propagation for data clustering.
\emph{Pattern Recognition} 157, 110933.

\bibitem{Yi2020AGSSNMF}
Yi, Y., Chen, Y., Wang, J., Lei, G., Dai, J., Zhang, H., 2020.
Joint feature representation and classification via adaptive graph semi-supervised nonnegative matrix factorization.
\emph{Signal Processing: Image Communication} 89, 115984.

\bibitem{Zhang2021MSNMF}
Zhang, K., Zhao, X., Peng, S., 2021.
Multiple graph regularized semi-supervised nonnegative matrix factorization with adaptive weights for clustering.
\emph{Engineering Applications of Artificial Intelligence} 106, 104499.

\bibitem{Li2025GOCNMF}
Li, J., Li, C., 2025.
Graph regularized one-hot-constrained nonnegative matrix factorization for data representation.
\emph{Engineering Applications of Artificial Intelligence} 156, 111224.

\bibitem{Zelnik2004SelfTuning}
Zelnik-Manor, L., Perona, P., 2004.
Self-tuning spectral clustering.
In: \emph{Advances in Neural Information Processing Systems 17}, pp.~1601-1608.

\bibitem{Zhu2003Harmonic}
Zhu, X., Ghahramani, Z., Lafferty, J., 2003.
Semi-supervised learning using Gaussian fields and harmonic functions.
In: \emph{Proceedings of the 20th International Conference on Machine Learning}, pp.~912-919.

\bibitem{Brier1950}
Brier, G.~W., 1950.
Verification of forecasts expressed in terms of probability.
\emph{Monthly Weather Review} 78~(1), 1-3.

\bibitem{Peng2021RSNMF}
Peng, S., Ser, W., Chen, B., Lin, Z., 2021.
Robust semi-supervised nonnegative matrix factorization for image clustering.
\emph{Pattern Recognition} 111, 107683.

\bibitem{Peng2022DCNMF}
Peng, S., Ser, W., Chen, B., Lin, Z., 2022.
Dual semi-supervised convex nonnegative matrix factorization for data representation.
\emph{Information Sciences} 585, 571-593.

\bibitem{Chavoshinejad2023S4NMF}
Chavoshinejad, J., Seyedi, S.~A., Tab, F.~A., Salahian, N., 2023.
Self-supervised semi-supervised nonnegative matrix factorization for data clustering.
\emph{Pattern Recognition} 137, 109282.

\bibitem{Li2025EDNMF}
Li, J., Shen, X., Li, C., Li, Y., 2025.
Elements discriminative non-negative matrix factorization for data clustering.
\emph{Engineering Applications of Artificial Intelligence} 156, 111210.

\bibitem{Qi2024GSSC}
Qi, T., Feng, X., Gao, B., Wang, K., 2024.
An end-to-end graph convolutional network for semi-supervised subspace clustering via label self-expressiveness.
\emph{Knowledge-Based Systems} 286, 111393.

\bibitem{Yin2023HSSNMF}
Yin, J., Peng, S., Yang, Z., Chen, B., Lin, Z., 2023.
Hypergraph based semi-supervised symmetric nonnegative matrix factorization for image clustering.
\emph{Pattern Recognition} 137, 109274.

\bibitem{Gan2018SaGSSL}
Gan, H., Li, Z., Wu, W., Luo, Z., Huang, R., 2018.
Safety-aware graph-based semi-supervised learning.
\emph{Expert Systems with Applications} 107, 243-254.

\bibitem{Wilcoxon1945}
Wilcoxon, F., 1945.
Individual comparisons by ranking methods.
\emph{Biometrics Bulletin} 1~(6), 80-83.

\bibitem{Song2013NEU}
Song, K., Yan, Y., 2013.
A noise robust method based on completed local binary patterns for hot-rolled steel strip surface defects.
\emph{Applied Surface Science} 285, 858-864.

\end{thebibliography}


\end{document}
